% Merged from three separately delivered research reports on the number of
% terms in the derivatives of a power tower.  See the README for provenance.
\documentclass[11pt,a4paper]{article}
\usepackage[T1]{fontenc}
\usepackage{lmodern}
\usepackage[margin=27mm]{geometry}
\usepackage{amsmath,amssymb,amsthm,mathtools}
\usepackage{booktabs,array,longtable}
\usepackage{microtype}
\usepackage{enumitem}
\usepackage{xcolor}
\usepackage{listings}
\usepackage{needspace}
\usepackage{fancyhdr}
\usepackage{hyperref}
\usepackage{bookmark}

\definecolor{linkink}{HTML}{294E47}
\hypersetup{colorlinks=true,linkcolor=linkink,citecolor=linkink,urlcolor=linkink,
 pdftitle={Term counts in the derivatives of power towers: OEIS A293239, A290268, A281434},
 pdfauthor={Research notes prepared for Vladimir Reshetnikov},
 pdfsubject={One framework, three sequences; none of the three headline questions is fully settled}}

\pagestyle{fancy}
\fancyhf{}
\fancyhead[L]{\small Term counts in the derivatives of power towers}
\fancyhead[R]{\small\leftmark}
\fancyfoot[C]{\thepage}
\renewcommand{\sectionmark}[1]{\markboth{#1}{}}
\setlength{\headheight}{14pt}
\setlength{\parskip}{4pt}
\setlength{\parindent}{0pt}
\setlength{\emergencystretch}{2em}

\numberwithin{equation}{section}
\newtheorem{theorem}{Theorem}[section]
\newtheorem{proposition}[theorem]{Proposition}
\newtheorem{lemma}[theorem]{Lemma}
\newtheorem{corollary}[theorem]{Corollary}
\theoremstyle{definition}
\newtheorem{definition}[theorem]{Definition}
\newtheorem{conjecture}[theorem]{Conjecture}
\theoremstyle{remark}
\newtheorem{remark}[theorem]{Remark}
\theoremstyle{plain}

\DeclareMathOperator{\supp}{supp}
\DeclareMathOperator{\Res}{Res}
\DeclareMathOperator{\sgn}{sgn}
\DeclareMathOperator{\atanh}{atanh}

\newcommand{\Z}{\mathbb Z}
\newcommand{\Q}{\mathbb Q}
\newcommand{\R}{\mathbb R}
\newcommand{\N}{\mathbb N}
\newcommand{\ZZ}{\mathcal Z}
\newcommand{\EE}{\mathcal E}
\newcommand{\U}{U}
\newcommand{\D}{\frac{d}{dx}}
\newcommand{\epsn}{\varepsilon_n}
\newcommand{\coef}[2]{[#1^{#2}]}
\newcommand{\coeff}[2]{[#1^{#2}]}
\newcommand{\ind}[1]{\mathbf{1}_{\{#1\}}}
% The three sources wrote the falling factorial two ways, (u)_{\underline n} and
% u^{\underline n}.  They are the same object -- each use in the text is
% accompanied by its explicit product expansion -- and the merge fixes one form.
\newcommand{\fall}[2]{#1^{\underline{#2}}}
\newcommand{\rise}[2]{#1^{\overline{#2}}}

\lstset{basicstyle=\ttfamily\small,columns=fullflexible,breaklines=true,
 frame=single,framerule=0.3pt,rulecolor=\color{black!25},keepspaces=true,
 showstringspaces=false,xleftmargin=5pt,xrightmargin=5pt,
 aboveskip=10pt,belowskip=10pt}

\title{\textbf{Term counts in the derivatives of power towers}\\[5pt]
\large One framework and three sequences:\\
OEIS A293239 \ ($x^x$), \ A290268 \ ($x^{x^2}$), \ A281434 \ ($x^{x^x}$)}
\author{Research notes prepared for Vladimir Reshetnikov}
\date{September 2026}

\begin{document}
\maketitle

\begin{abstract}
For a function $f$, let $a(n)$ be the number of distinct monomials in the fully
collected $n$th derivative $f^{(n)}$.  This article treats three instances of
that question at once: $f=x^x$ (OEIS A293239), $f=x^{x^2}$ (A290268) and
$f=x^{x^x}$ (A281434).

Part~\ref{fw:sec} proves, once, the apparatus the three share: a canonical
differential-polynomial form $f^{(n)}=f\cdot P_n$ or $f\cdot x^{-n}P_n$ with a
first-order recurrence for $P_n$ over $\Z$, integrality of its coefficients, and
the envelope bound that turns an upper estimate into a cancellation count.  It
also proves the linear independence of logarithmic monomials that makes ``the
number of terms'' an invariant of $f$ rather than of a differentiation
algorithm, and the resulting identity $a(n)=|\supp P_n|$ --- for two of the
three frames; $x^{x^x}$ falls outside that statement and supplies its own, as
Part~\ref{fw:sec} is careful to say.  Two normalizations are stated, because neither one covers
all three parts in the frames those parts use; Remark~\ref{fw:rem:twoforms} is
careful about what that does and does not mean.

The three remaining parts are the three original investigations, unabridged.
They reach different answers by different routes.  \textbf{None of the three
headline questions is fully settled here, and this article does not obscure
that.}  Two of the three entries print a conjectured closed form and neither is
proved; the third prints no closed form, and there the growth \emph{order} is
settled while the constant is not.  For $x^x$ the
term count is reduced exactly to the zeros of the first-kind Lehmer--Comtet
triangle A008296, giving $a(n)=1+\binom{n+1}{2}-Z(n)$ and $a(n)=\Theta(n^2)$,
but the principal closed formula needs a nonvanishing statement that is isolated
and left open; separately, the recurrence printed in the entry with the range
$n>6$ is \emph{disproved} at $n=8$.  For $x^{x^2}$ the conjectured expression is
proved to be an upper bound, every predicted cancellation is explained, and
$a(n)=\Theta(n^2)$ is established, but the conjectured leading constant is not.
For $x^{x^x}$ two explicit polynomial bounds give $a(n)=\Theta(n^3)$
unconditionally, while the sharper equivalent $a(n)\sim\tfrac23n^3$ is not
proved --- and there, unlike the other two, the obstruction is not a
nonvanishing statement at all but a density one, because additional zeros
provably exist.  Part~\ref{syn:sec} sets the three side by side, says what the
shared method does and does not buy, and records the one piece of structure they
share beyond it: every cancellation any of them can prove is a parity zero of
the same falling factorial.

All computations reported here are exact.  Every modular certificate is
confirmed by exact integer arithmetic before a coefficient is declared to
vanish, so no claim rests on a probabilistic zero test; but every computation
certifies a finite range and none of them is an all-$n$ proof.
\end{abstract}

\setcounter{tocdepth}{2}
\tableofcontents
\clearpage

\section*{How this article is arranged}
\addcontentsline{toc}{section}{How this article is arranged}

Three separate reports were written on three OEIS sequences that ask the same
question about three different functions.  They are collected here as one
article because they share a construction, and kept as four distinct parts
because the questions they answer are different: three functions, three
recurrences, three answers.  They do share one thing beyond the construction,
and Section~\ref{syn:share} sets it out --- all three reduce to integer zeros of
derivatives and differences of the same falling factorial --- but no argument in
any part carries over to another.  Nothing that was proved in any of the three
has been dropped, weakened or strengthened in the merge.

\begin{itemize}[leftmargin=1.6em,itemsep=2pt]
\item \textbf{Part~\ref{fw:sec}} is new to this article.  It states and proves
 the machinery that the three reports had each established separately, and is
 explicit about where that generality stops: the canonical form, the recurrence
 and the envelope bound cover all three, the independence statement covers two
 of the three.  It settles none of the
 three counting problems and is not intended to: the content of the subject is
 in the nonvanishing questions, and those do not transfer between the three
 functions.
\item \textbf{Parts~\ref{xx:part}, \ref{xxb:part} and~\ref{xxc:part}} are the
 three original articles.  Each opens with a status box saying what it does and
 does not prove, and a note on how it instantiates Part~\ref{fw:sec}.  Each part
 retains the setup section it was written with, so it can still be read on its
 own; those sections are now instances of Part~\ref{fw:sec} rather than
 independent developments, and are flagged as such where they begin.
\item \textbf{Part~\ref{syn:sec}} compares the three.
\end{itemize}

The parts keep the notation of their sources, which is not uniform across them.
Section~\ref{fw:notation} sets out what a reader crossing between parts needs.
Two things are worth carrying in advance: \textbf{$j$ is the exponent of
$\log x$ in Parts~\ref{xx:part} and~\ref{xxb:part} but the exponent of $x^{-1}$
in Part~\ref{xxc:part}}, where $\ell$ is the logarithm exponent; and
\textbf{$z$ is $x^2$ in Part~\ref{xxb:part}, $\log x$ in Part~\ref{xxc:part},
and a formal generating-function variable in Part~\ref{xx:part}}.  Also
$a(n)$ means a different sequence in each of the three parts.

Each part ends with a description of the code and data that accompanied it.
Those descriptions have been updated to the merged package, in which every
script lives under \texttt{code/} and every generated table under
\texttt{data/}, and the article is \texttt{article.tex} and
\texttt{article.pdf}.  The three groups of scripts do not collide and each was
run from the merged layout before this article was built.

These are AI-assisted research notes.  None of the three has been refereed or
machine-checked.

\section{The common framework}\label{fw:sec}

Each of the three parts that follow studies a different function, obtains a
different recurrence and arrives at a different answer.  What they share is the
apparatus in this part: a canonical polynomial form for the $n$th derivative, a
uniqueness statement that makes ``the number of terms'' well defined, a
first-order recurrence on that form, and the reduction of the counting problem
to a nonvanishing problem.  Each of the three source reports established this
apparatus for its own function.  Here it is proved once.

Nothing in this part decides any of the three counting problems.  It is
deliberately cheap: every statement below is an exercise, and the work in this
subject is entirely in the nonvanishing questions that Parts~\ref{xx:part},
\ref{xxb:part} and~\ref{xxc:part} take up.

\subsection{Setting}\label{fw:setting}

Throughout, functions are real-valued and smooth on $(0,\infty)$, and $f$ is
positive there, so $\log f$ is defined and $f'/f$ makes sense.

\begin{definition}[Differential frame]\label{fw:def:frame}
A \emph{differential frame} is a tuple $w=(w_1,\dots,w_m)$ of smooth functions
on $(0,\infty)$ together with the polynomial ring $R=\Z[w_1,\dots,w_m]$, where
the $w_i$ are treated as independent indeterminates.  We write $\supp P$ for the
set of exponent vectors of the monomials occurring in $P\in R$ with a nonzero
coefficient, and $|\supp P|$ for its cardinality.
\end{definition}

The distinction between $R$ as an abstract polynomial ring and the functions the
$w_i$ denote is the point of Section~\ref{fw:independence}: a priori a
polynomial identity in $R$ is a stronger statement than the corresponding
identity of functions, and it is the converse that has to be proved.

\subsection{The canonical form and its recurrence}\label{fw:canonical}

\begin{theorem}[Canonical differential-polynomial form]\label{fw:thm:canonical}
Let $w$ be a differential frame with ring $R=\Z[w_1,\dots,w_m]$, and let
$D\colon R\to R$ be the $\Z$-derivation determined by $Dw_i=\delta_i$ for some
$\delta_1,\dots,\delta_m\in R$.
\begin{enumerate}[label=\textup{(\alph*)}]
\item\label{fw:can:a}
 \textup{(Unnormalized form.)}  Suppose
 \[
  w_i'=\delta_i(w)\quad(1\le i\le m),
  \qquad \frac{f'}{f}=M(w)\ \text{ for some }M\in R .
 \]
 Then for every $n\ge0$,
 \begin{equation}\label{fw:eq:cana}
  f^{(n)}=f\cdot P_n(w),\qquad
  P_0=1,\qquad P_{n+1}=M\,P_n+D\,P_n ,
 \end{equation}
 and $P_n\in R$; in particular every coefficient of $P_n$ is an integer.
\item\label{fw:can:b}
 \textup{(Form normalized by $x^{-n}$.)}  Suppose instead
 \[
  w_i'=\frac1x\,\delta_i(w)\quad(1\le i\le m),
  \qquad x\,\frac{f'}{f}=M(w)\ \text{ for some }M\in R .
 \]
 Then for every $n\ge0$,
 \begin{equation}\label{fw:eq:canb}
  f^{(n)}=f\cdot x^{-n}P_n(w),\qquad
  P_0=1,\qquad P_{n+1}=(M-n)\,P_n+D\,P_n ,
 \end{equation}
 and again $P_n\in R$.
\end{enumerate}
\end{theorem}

\begin{proof}
For~\ref{fw:can:a}, induct on $n$.  The case $n=0$ is $P_0=1$.  Assuming
$f^{(n)}=fP_n(w)$ and differentiating,
\[
 f^{(n+1)}
 =f'P_n(w)+f\sum_{i=1}^m w_i'\,\frac{\partial P_n}{\partial w_i}(w)
 =f\Bigl(M(w)P_n(w)+\sum_{i=1}^m\delta_i(w)\frac{\partial P_n}{\partial w_i}(w)\Bigr),
\]
which is $f\bigl(MP_n+DP_n\bigr)(w)$.  Since $M,\delta_i\in R$ and $P_0=1\in R$,
induction also gives $P_n\in R$, and $R$ consists of polynomials with integer
coefficients.

For~\ref{fw:can:b}, the case $n=0$ reads $f=f\,x^{0}\cdot1$.  Assume
$f^{(n)}=fx^{-n}P_n(w)$ and differentiate:
\[
 f^{(n+1)}
 =f'x^{-n}P_n-nfx^{-n-1}P_n+fx^{-n}\sum_i w_i'\frac{\partial P_n}{\partial w_i}.
\]
Multiplying by $x^{n+1}/f$ and using $xf'/f=M$ and $xw_i'=\delta_i$ gives
\[
 \frac{x^{n+1}f^{(n+1)}}{f}=\bigl(M-n\bigr)P_n+\sum_i\delta_i\frac{\partial P_n}{\partial w_i}
 =\bigl((M-n)P_n+DP_n\bigr)(w),
\]
which is the assertion for $n+1$.  Integrality follows as before.
\end{proof}

\begin{remark}[Why both clauses are kept]\label{fw:rem:twoforms}
Neither clause covers all three parts \emph{in the frames those parts use}, and
since the parts are reproduced here unchanged, both clauses are needed.

Part~\ref{xxb:part} uses the frame $z=x^2$, $L=\log x$, for which
$z'=2x\notin\Z[z,L]$: clause~\ref{fw:can:a} does not apply, and
clause~\ref{fw:can:b} applies verbatim.  Conversely Part~\ref{xxc:part} uses
$u=x^x$, $v=x^{-1}$, $z=\log x$, for which $xu'=u(z+1)/v\notin\Z[u,v,z]$:
clause~\ref{fw:can:b} does not apply, while clause~\ref{fw:can:a} does.

It would overstate the situation to say that no single clause \emph{could} have
covered all three.  Clause~\ref{fw:can:b} in fact does, once the frames are
changed so that $x$ appears in place of $x^{-1}$:
\[
 x^x:\ (x,\log x),\qquad x^{x^2}:\ (x^2,\log x),\qquad x^{x^x}:\ (x^x,x,\log x),
\]
each of which satisfies $xw_i'\in R$ and $xf'/f\in R$, and each of which gives
the same term counts.  For $x^{x^x}$ this is legitimate only because the
exponent of $x^{-1}$ in that part's envelope never exceeds $n$, so multiplying
by $x^{n}$ really does clear it.  What one cannot do is change the frames and
keep the parts: each part's transition table, envelope, coefficient indices and
data files are written in its own variables, and rewriting some two and a half
thousand lines of dense manipulation to gain a cosmetic uniformity in this
section would be a poor trade.  So the parts keep their frames and this theorem
keeps two clauses.

One genuine obstruction is worth recording separately, since it is about
independence rather than convenience.  In Part~\ref{xxb:part} one might try to
make clause~\ref{fw:can:a} apply by adjoining $y=x^{-1}$ to the frame
$(x^2,\log x)$.  The derivatives then are polynomial, but $zy^2=1$ identically,
so the monomials of the enlarged frame are no longer independent as functions
and $|\supp P_n|$ stops counting anything at all.  That failure is not a matter
of taste.
\end{remark}

Each of the three parts instantiates Theorem~\ref{fw:thm:canonical} in one line;
the resulting recurrences are collected in Table~\ref{fw:tab:instances}.

\begin{table}[htbp]
\centering
\small
\begin{tabular}{@{}llll@{}}
\toprule
Part & $f$ & frame and clause & recurrence for $P_{n+1}$\\
\midrule
\ref{xx:part} & $x^{x}$ &
 $t=x^{-1},\ L=\log x$; \ref{fw:can:a} &
 $(1+L)P_n+t\,\partial_LP_n-t^2\,\partial_tP_n$\\[2pt]
\ref{xxb:part} & $x^{x^2}$ &
 $z=x^2,\ L=\log x$; \ref{fw:can:b} &
 $\bigl(z(2L+1)-n\bigr)P_n+2z\,\partial_zP_n+\partial_LP_n$\\[2pt]
\ref{xxc:part} & $x^{x^x}$ &
 $u=x^x,\ v=x^{-1},\ z=\log x$; \ref{fw:can:a} &
 $u(z^2{+}z{+}v)P_n+u(z{+}1)\partial_uP_n-v^2\partial_vP_n+v\,\partial_zP_n$\\
\bottomrule
\end{tabular}
\caption{The three instances.  In each case the displayed recurrence is exactly
the one derived independently in the corresponding part
(Propositions~\ref{xx:prop:Prec}, \ref{xxb:prop:recurrence}
and~\ref{xxc:prop:recurrence}).  Note that $z$ denotes $x^2$ in
Part~\ref{xxb:part} and $\log x$ in Part~\ref{xxc:part}; see
Section~\ref{fw:notation}.}
\label{fw:tab:instances}
\end{table}

\subsection{Independence, and why the count is well posed}\label{fw:independence}

Counting ``the terms of the $n$th derivative'' is meaningless until one fixes
what a term is, because a symbolic differentiator can be made to produce almost
any number of summands by factoring differently.  The following classical
statement is what makes the count an invariant of the function rather than of an
algorithm.  Parts~\ref{xx:part} and~\ref{xxb:part} each prove the special case
they need by this argument.  Part~\ref{xxc:part}'s frame is not of this form,
for the reason given after Corollary~\ref{fw:cor:welldefined}, and it proves its
own independence statement by a related but genuinely different asymptotic
argument.

\begin{theorem}[Independence of logarithmic monomials]\label{fw:thm:independence}
Let $(\lambda_1,j_1),\dots,(\lambda_N,j_N)$ be distinct pairs with
$\lambda_k\in\R$ and $j_k\in\Z_{\ge0}$.  Then the functions
$x^{\lambda_k}(\log x)^{j_k}$ are linearly independent over $\R$ on every
interval $(a,\infty)$ with $a\ge0$.
\end{theorem}

\begin{proof}
Suppose $\sum_k c_kx^{\lambda_k}(\log x)^{j_k}=0$ on $(a,\infty)$ with not all
$c_k$ zero.  Substitute $x=e^s$ and group by exponent: the identity becomes
\begin{equation}\label{fw:eq:indep}
 \sum_{\lambda}e^{\lambda s}p_\lambda(s)=0\qquad(s>\log\max(a,1)),
\end{equation}
where $\lambda$ runs over the distinct values among the $\lambda_k$ and
$p_\lambda(s)=\sum_{k:\lambda_k=\lambda}c_ks^{j_k}$ is a polynomial.  Because the
pairs are distinct, $p_\lambda$ is not the zero polynomial for at least one
$\lambda$.  Let $\Lambda$ be the largest such $\lambda$.  Dividing
\eqref{fw:eq:indep} by $e^{\Lambda s}$ gives
\[
 p_\Lambda(s)=-\sum_{\lambda<\Lambda}e^{(\lambda-\Lambda)s}p_\lambda(s)
 \xrightarrow[s\to+\infty]{}0,
\]
since each exponent $\lambda-\Lambda$ is strictly negative and an exponential
decay beats a polynomial.  No nonzero polynomial tends to $0$ at $+\infty$, so
$p_\Lambda=0$, a contradiction.
\end{proof}

\begin{corollary}[The count is well defined]\label{fw:cor:welldefined}
Let $f$ and $w$ be as in either clause of Theorem~\textup{\ref{fw:thm:canonical}},
and suppose that the monomials $w^\alpha$ of the frame correspond to pairwise
distinct pairs $(\lambda(\alpha),j(\alpha))$ in the sense that
$w^\alpha=x^{\lambda(\alpha)}(\log x)^{j(\alpha)}$.  Then the representation
of $f^{(n)}$ in \eqref{fw:eq:cana}, respectively \eqref{fw:eq:canb}, is unique,
and the number of distinct monomials in the fully collected $n$th derivative is
\begin{equation}\label{fw:eq:count}
 a(n)=|\supp P_n| .
\end{equation}
\end{corollary}

\begin{proof}
Two representations would differ by a finite vanishing linear combination of the
functions $w^\alpha$, hence of pairwise distinct $x^{\lambda}(\log x)^{j}$,
which Theorem~\ref{fw:thm:independence} forbids.  Uniqueness turns
\eqref{fw:eq:count} into a definition that does not refer to any particular way
of performing the differentiation.
\end{proof}

The hypothesis of Corollary~\ref{fw:cor:welldefined} is verified separately in
each part, since it depends on the frame: for $x^x$ the monomials $t^kL^j$ are
$x^{-k}(\log x)^j$; for $x^{x^2}$ the monomials $z^kL^j$ carry the extra factor
$x^{-n}$ and give $x^{2k-n}(\log x)^j$; for $x^{x^x}$ the variable $u=x^x$ is
itself not of the form $x^\lambda$ with constant $\lambda$, and
Part~\ref{xxc:part} accordingly proves its independence statement by a related
but genuinely different asymptotic argument, which is retained there in full.

\subsection{Envelopes, upper bounds and exact computation}\label{fw:envelope}

\begin{proposition}[Envelope bound]\label{fw:prop:envelope}
Write $P_{n+1}=\Phi_n(P_n)$ for either recurrence of
Theorem~\textup{\ref{fw:thm:canonical}}, so that $\Phi_n(P)=MP+DP$ in
clause~\textup{\ref{fw:can:a}} and $\Phi_n(P)=(M-n)P+DP$ in
clause~\textup{\ref{fw:can:b}}; note that in the second case the operator
depends on $n$, although $\supp\bigl((M-n)P\bigr)\subseteq\supp(MP)\cup\supp P$
for every $n$, so the support behaviour does not.  Suppose $\mathcal
E_0,\mathcal E_1,\dots$ are finite sets of exponent vectors with $\supp
P_0\subseteq\mathcal E_0$ and with the property that $\Phi_n$ maps any
polynomial supported in $\mathcal E_n$ to a polynomial supported in $\mathcal
E_{n+1}$.  Then $\supp P_n\subseteq\mathcal E_n$ for every $n$; and if moreover
the count is well defined in the sense of
Corollary~\textup{\ref{fw:cor:welldefined}}, or by the part's own substitute for
it, then
\[
 a(n)\le|\mathcal E_n| .
\]
\end{proposition}

\begin{proof}
Immediate induction for the support statement.  Only the direction matters:
$\Phi_n$ sends each monomial to a bounded set of monomials, so a cell outside
$\mathcal E_{n+1}$ receives no contribution at all.  Cancellation between
contributions can remove a cell of $\mathcal E_{n+1}$ from the support, but it
can never create one outside.  The bound on $a(n)$ is then
\eqref{fw:eq:count}.
\end{proof}

The deficit
\begin{equation}\label{fw:eq:deficit}
 \Delta(n)=|\mathcal E_n|-a(n)\ \ge 0
\end{equation}
is therefore exactly the number of cells inside the envelope whose coefficient
vanishes.  This is the shape every one of the three problems takes: an envelope
that is easy to count, and a cancellation count that is not.  The three parts
differ precisely in how much they can say about $\Delta$.

\begin{remark}[The parts count deficits differently]\label{fw:rem:deficits}
A reader moving between the parts must not compare their deficit numbers
directly.  Part~\ref{xxc:part} writes $\Delta(n)$ for \emph{all} vanishing cells
of its envelope, including the ones it proves must vanish; that quantity is
itself a published sequence, OEIS A352697.  Part~\ref{xx:part} writes $E(n)$
for only the zeros \emph{outside} the families it has already proved to vanish,
so $E$ is the residual defect and is conjectured to be identically zero, which
$\Delta$ certainly is not.  Part~\ref{xxb:part} likewise subtracts its proved
reflection zeros before counting.  Both conventions are useful and each part
keeps its own; the generic $\Delta$ of \eqref{fw:eq:deficit} is the first kind.
\end{remark}

\begin{proposition}[Exact and modular computation]\label{fw:prop:computation}
With notation as above, $P_0,\dots,P_N$ can be computed in
$O\bigl(\sum_{n\le N}|\mathcal E_n|\bigr)$ integer ring operations, using only
integer arithmetic; the implied constant is the frame's transition count
$|\supp M|+\sum_i|\supp\delta_i|$.  Moreover, for any modulus $q\ge2$, a
coefficient of $P_n$ whose residue modulo $q$ is nonzero is nonzero.
\end{proposition}

\begin{proof}
Each recurrence applies at most $|\supp M|+\sum_i|\supp\delta_i|$
monomial-to-monomial transitions to each cell of the support --- a constant
depending on the frame and not on $n$ --- so one step costs $O(|\mathcal E_n|)$
operations; the coefficients are integers by
Theorem~\ref{fw:thm:canonical}.  Reduction modulo $q$ is a ring homomorphism, so
it commutes with the recurrence; a nonzero image forces a nonzero preimage.
\end{proof}

Proposition~\ref{fw:prop:computation} is the reason the computations reported in
the three parts are certifications and not samples: a nonzero residue settles
nonvanishing outright.  The converse direction is the weak one --- a zero
residue on its own certifies nothing --- and the three parts close it in three
different ways, which is worth stating precisely rather than in one sentence:

\begin{itemize}[leftmargin=1.6em,itemsep=2pt]
\item Part~\ref{xx:part} computes the triangle in exact integer arithmetic
 outright, so no residue has to be resolved; its modular lemma is used for a
 different purpose, to rule out integer roots of a fixed polynomial over a whole
 diagonal.
\item Part~\ref{xxb:part} runs two moduli and checks that the two lists of
 vanishing residues are disjoint, so at least one residue of every coefficient
 is nonzero.  Its exact integer run stops at $n=200$; the certification to
 $3000$ is modular.
\item Part~\ref{xxc:part} re-evaluates every vanishing residue with an
 independent exact coefficient formula, and never accepts a zero residue as
 evidence.
\end{itemize}

\noindent
All three are deterministic.  No statement anywhere in this article rests on a
probabilistic zero test.  None of them, however, is an all-$n$ proof: each
certifies a finite range, and each part says so where it reports the range.

\subsection{Notation across the parts}\label{fw:notation}

The three parts below retain the notation of the reports they come from.  That
notation is not consistent between them, and rewriting it would have meant
altering some two and a half thousand lines of dense manipulation with no way to
check the result; declaring the clash is safer than silently repairing it.

The table that matters is not a dictionary of letters --- nearly every letter is
reused somewhere as a summation index, a residue or a generating-function
variable, and no short table can be both complete and correct about that.  What
a reader does need, before crossing from one part to another, is the canonical
monomial each part counts, and in particular which of its indices is the
exponent of $\log x$.

\begin{center}
\begin{tabular}{@{}llll@{}}
\toprule
 & Part~\ref{xx:part}: $x^x$ & Part~\ref{xxb:part}: $x^{x^2}$ &
 Part~\ref{xxc:part}: $x^{x^x}$\\
\midrule
frame & $t=x^{-1}$, $L=\log x$ & $z=x^2$, $L=\log x$ &
 $u=x^x$, $v=x^{-1}$, $z=\log x$\\[3pt]
monomial & $x^{\,x-k}(\log x)^{j}$ & $x^{\,x^2+2k-n}(\log x)^{j}$ &
 $x^{\,x^x+kx-j}(\log x)^{\ell}$\\[3pt]
$\log x$ exponent & $j$ & $j$ & $\boldsymbol{\ell}$\\[3pt]
power of $x^{-1}$ & $k$ & --- (absorbed in $x^{-n}$) & $\boldsymbol{j}$\\[3pt]
normalization & $f^{(n)}=x^xP_n$ & $f^{(n)}=f\,x^{-n}P_n$ & $f^{(n)}=f\,P_n$\\
\bottomrule
\end{tabular}
\end{center}

\noindent
Two warnings follow from that table, and they are the two places where a reader
who is skimming will go wrong.

\begin{enumerate}[leftmargin=1.8em,itemsep=3pt]
\item \textbf{$j$ changes role.}  In Parts~\ref{xx:part} and~\ref{xxb:part} the
 index $j$ is the exponent of $\log x$.  In Part~\ref{xxc:part} it is the
 exponent of $x^{-1}$, and the exponent of $\log x$ is $\ell$.  Every envelope,
 every bound and every cancellation family in Part~\ref{xxc:part} is indexed in
 the second convention.
\item \textbf{$z$ has three meanings.}  It is $x^2$ in Part~\ref{xxb:part} and
 $\log x$ in Part~\ref{xxc:part} --- a direct contradiction --- and in
 Part~\ref{xx:part} it is neither: there it is the formal variable in which the
 Lehmer--Comtet numbers are defined, through $\Phi(z)=(1+z)\log(1+z)$ and
 $b(m,r)=\tfrac{m!}{r!}[z^m]\Phi(z)^r$.
\end{enumerate}

\noindent
Beyond these, the letters $u$, $v$, $t$, $q$, $U$, $Z$, $E$ and $\Delta$ all
occur in more than one part with unrelated meanings, as frame variables in one
place and as summation indices, residues, primes, proposed closed forms or
generating-function variables in another.  Each is introduced where it is used,
and within any one paragraph the meaning is fixed; but nothing in this article
guarantees that a letter carries the same meaning two parts later, and the
reader should not assume it.  Finally, $a(n)$ denotes a different sequence in
each of the three parts --- its own --- and the deficit conventions differ as
described in Remark~\ref{fw:rem:deficits}.

\clearpage
\section{\texorpdfstring{$x^x$}{x^x}: OEIS A293239}\label{xx:part}

\begin{quote}\small
\textbf{Framework.}  This part instantiates
Theorem~\ref{fw:thm:canonical}\ref{fw:can:a} with the frame $t=x^{-1}$,
$L=\log x$; Proposition~\ref{xx:prop:Prec} below is that instance, derived here
independently.  Lemma~\ref{xx:lem:independence} is the special case of
Theorem~\ref{fw:thm:independence} this part needs.  Both are retained in full so
that the part can be read on its own.  Throughout this part $a(n)$ denotes the
term count for $x^x$.

\medskip
\textbf{Abstract of the original report.}\itshape
The sequence A293239 counts the nonzero terms in the fully expanded $n$th
derivative of $x^x$. Its proposed generating function is rational, and its
proposed eventual formula is a quadratic quasipolynomial of period four.
We reduce the counting problem exactly to the zero set of the first-kind
Lehmer--Comtet numbers. All cancellations predicted by the conjecture are
proved: a symmetry family and one additional zero. The converse assertion,
excluding every further zero, is not established here; consequently this
report does \upshape\emph{not}\itshape{} claim a complete proof of the principal
closed formula.
The unconditional results include an exact defect formula, upper and lower
bounds proving quadratic-order growth, nonvanishing on every prime-offset
diagonal with an exact prime-adic valuation, and complete zero classifications
for the first three columns and offsets one through sixteen. Independently
implemented exact computations confirm the formula through $n=3000$.
A separate conjectured recurrence in the OEIS entry is false with its printed
range: it already fails at $n=8$. For the proposed sequence its correct
eventual starting point is $n=14$. The accompanying code and certificates
distinguish finite verification from statements proved for infinitely many
indices.
\end{quote}

\subsection{The question and the status of the answer}

For $x>0$, put $f(x)=x^x=\exp(x\log x)$. Let $a(n)$ be the number of nonzero
summands after $f^{(n)}(x)$ is expanded into the monomials
\begin{equation}\label{xx:eq:monomials}
 x^{x-k}(\log x)^j,\qquad k,j\in\Z_{\ge0},
\end{equation}
and coefficients of identical monomials are collected. In particular, an
unexpanded expression such as $x^x(1+\log x)^n$ is not one term for this
purpose. This is the convention in OEIS A293239, introduced by Vladimir
Reshetnikov in 2017 \cite{oeis293239}.

The entry attributes the rational generating function and closed formula to
Colin Barker and the proposed asymptotic equivalent to V\'aclav
Kote\v{s}ovec. To separate a proved object from the proposed answer, define
\begin{equation}\label{xx:eq:q}
 q(0)=1,\qquad
 q(n)=1+\frac{n(n+1)}2-\left\lfloor\frac{n-1}{4}\right\rfloor
             -\ind{n\ge8}\quad(n\ge1).
\end{equation}
The principal conjecture is $a(n)=q(n)$. For $n\ge8$, this says simply
\begin{equation}\label{xx:eq:eventual}
 a(n)\stackrel{?}{=}\frac{n(n+1)}2-\left\lfloor\frac{n-1}{4}\right\rfloor.
\end{equation}
The proposed generating function is
\begin{equation}\label{xx:eq:targetgf}
 Q(z):=\sum_{n\ge0}q(n)z^n
 =\frac{1+z^2+z^3+z^6-z^8+z^9+z^{12}-z^{13}}
 {(1-z)^2(1-z^4)}.
\end{equation}
The rational function in \eqref{xx:eq:targetgf} is indeed the generating function
of the explicitly defined sequence $q$. What remains to be proved is its
identification with the derivative count $a$.

\subsubsection{What is proved}
The central unconditional identity will be
\begin{equation}\label{xx:eq:introdefect}
 a(n)=q(n)-E(n),
\end{equation}
where $E(n)$ counts a precisely defined set of additional zero coefficients.
In particular, $E(n)$ is a nonnegative, nondecreasing integer. Thus the
proposed answer is an upper bound, not a formula that can be exceeded by a
previously overlooked cancellation.

We prove, for all $n\ge0$,
\begin{equation}\label{xx:eq:introbounds}
 n+1+\left\lfloor\frac{n^2}{4}\right\rfloor\le a(n)\le q(n).
\end{equation}
Consequently $a(n)=\Theta(n^2)$. This does not establish the sharper claim
$a(n)\sim n^2/2$; that claim is equivalent to $E(n)=o(n^2)$.
We also prove infinite nonvanishing results that substantially restrict where
additional zeros could lie. The remaining gap is stated explicitly in
Conjecture~\ref{xx:conj:zeros}.

\subsubsection{A correction that is unconditional}
The OEIS entry also proposes
\begin{equation}\label{xx:eq:recprinted}
 a(n)=2a(n-1)-a(n-2)+a(n-4)-2a(n-5)+a(n-6)\qquad(n>6).
\end{equation}
This range is incorrect. Exact low-order differentiation gives
\[
 a(2)=4,\quad a(3)=7,\quad a(4)=11,\quad a(6)=21,\quad
 a(7)=28,\quad a(8)=35,
\]
whereas the right side of \eqref{xx:eq:recprinted} at $n=8$ is
\[
 2\cdot28-21+11-2\cdot7+4=36\ne35.
\]
For $q$, the recurrence holds for every $n\ge14$, and it fails at $n=13$.
This correction does not disprove the proposed generating function or
closed formula. Section~\ref{xx:sec:gf} explains the distinction.

\subsection{A canonical polynomial for each derivative}

Set $t=x^{-1}$ and $L=\log x$. There is a polynomial $P_n\in\Z[t,L]$ such that
\begin{equation}\label{xx:eq:Pdefinition}
 f^{(n)}(x)=x^x P_n(x^{-1},\log x).
\end{equation}
Differentiating this expression and using $t'=-t^2$, $L'=t$, and
$(x^x)'=x^x(1+L)$ gives the following exact recurrence.

\begin{proposition}\label{xx:prop:Prec}
The derivative polynomials satisfy
\begin{equation}\label{xx:eq:Prec}
 P_0=1,\qquad P_{n+1}=(1+L)P_n+t\frac{\partial P_n}{\partial L}
                         -t^2\frac{\partial P_n}{\partial t}.
\end{equation}
In particular, they have integer coefficients and can be computed using
integer arithmetic only.
\end{proposition}

For example,
\begin{align*}
 P_1&=1+L,\\
 P_2&=1+2L+L^2+t,\\
 P_3&=1+3L+3L^2+L^3+3t+3tL-t^2.
\end{align*}
Thus $a(3)=7$, consistently with the example in the entry.

\begin{lemma}\label{xx:lem:independence}
The monomials in \eqref{xx:eq:monomials} are linearly independent over $\R$
when only finitely many are used. Consequently the collected expansion is
unique, and $a(n)=|\supp P_n|$.
\end{lemma}
\begin{proof}
Divide a putative relation by the positive function $x^x$ and set $x=e^s$.
The result is $\sum_{k=0}^K e^{-ks}p_k(s)=0$, with polynomial $p_k$.
If some $p_k$ is nonzero, let $k_0$ be the smallest such index. Multiplication
by $e^{k_0s}$ shows that $p_{k_0}(s)$ tends to zero as $s\to+\infty$, because
all remaining summands decay exponentially times a polynomial. No nonzero
polynomial has that property. This is a contradiction.
\end{proof}

The independence statement matters: the variables $t$ and $L$ are related
by $t=e^{-L}$ after substitution, but this relation causes no additional
finite polynomial identities.

\subsection{The Lehmer--Comtet coefficient triangle}

\subsubsection{Definition and generating function}
Put
\begin{equation}\label{xx:eq:phi}
 \Phi(z)=(1+z)\log(1+z)
 =z+\sum_{d\ge1}\frac{(-1)^{d-1}}{d(d+1)}z^{d+1}.
\end{equation}
Define
\begin{equation}\label{xx:eq:bdef}
 b(m,r)=\frac{m!}{r!}\coef{z}{m}\Phi(z)^r
 \quad(m,r\ge0).
\end{equation}
Here $b(0,0)=1$, $b(m,0)=0$ for $m>0$, and $b(m,r)=0$ when $r>m$.
These are the first-kind Lehmer--Comtet numbers, tabulated in OEIS A008296
\cite{oeis8296}. The general derivative problem and these numbers belong to
the classical literature of Comtet, Lehmer, and Gould
\cite{comtet,lehmer,gould}. A recent treatment of their identities is
Jeong's preprint \cite{jeong}. We derive the identities needed here rather
than assuming a nonvanishing conclusion from those formulas.

\begin{theorem}[Exact coefficient formula]\label{xx:thm:expansion}
For every $n\ge0$,
\begin{equation}\label{xx:eq:expansion}
 \boxed{\displaystyle
 f^{(n)}(x)=x^x\sum_{m=0}^n\binom nm (\log x)^{n-m}
                         \sum_{r=0}^m b(m,r)x^{r-m}.}
\end{equation}
Equivalently,
\begin{equation}\label{xx:eq:cformula}
 \coef{t}{k}\coef{L}{j}P_n(t,L)
 =\binom nj\,b(n-j,n-j-k),
\end{equation}
with out-of-range entries interpreted as zero.
\end{theorem}
\begin{proof}
Taylor expansion at a fixed $x>0$ gives
\[
 \sum_{n\ge0}\frac{f^{(n)}(x)}{x^x}\frac{u^n}{n!}
 =\frac{(x+u)^{x+u}}{x^x}
 =\exp(u\log x)\exp\bigl(x\Phi(u/x)\bigr).
\]
Using \eqref{xx:eq:bdef}, expand the second exponential:
\begin{align*}
 \exp\bigl(x\Phi(u/x)\bigr)
 &=\sum_{r\ge0}\frac{x^r}{r!}\Phi(u/x)^r\\
 &=\sum_{m\ge0}\frac{u^m}{m!}\sum_{r=0}^m b(m,r)x^{r-m}.
\end{align*}
Multiplication by $\exp(u\log x)$ and extraction of the coefficient of
$u^n$ proves \eqref{xx:eq:expansion}. All manipulations can also be performed
formally, since only finitely many terms affect any prescribed coefficient.
\end{proof}

\subsubsection{A shifted falling-factorial formula}
Let
\[
 F_m(y)=\fall{y}{m}=\prod_{h=0}^{m-1}(y-h),\qquad F_0(y)=1.
\]

\begin{proposition}\label{xx:prop:falling}
For $0\le r\le m$,
\begin{equation}\label{xx:eq:falling}
 b(m,r)=\frac{F_m^{(r)}(r)}{r!}
       =\coef{u}{r}\prod_{h=0}^{m-1}(u+r-h).
\end{equation}
In particular, all $b(m,r)$ are integers.
\end{proposition}
\begin{proof}
The binomial-series identity is
$(1+z)^y=\sum_{m\ge0}F_m(y)z^m/m!$.
Differentiate $r$ times in $y$ and set $y=r$. The left side becomes
$(1+z)^r\log(1+z)^r=\Phi(z)^r$. Comparing coefficients gives the first
formula. The second is the polynomial Taylor formula at $y=r$.
\end{proof}

Writing $F_m(y)=\sum_{\ell=0}^m s(m,\ell)y^\ell$ in terms of the signed
Stirling numbers of the first kind also gives
\begin{equation}\label{xx:eq:stirling}
 b(m,r)=\sum_{\ell=r}^m\binom\ell r r^{\ell-r}s(m,\ell).
\end{equation}
This is another exact formula, but its alternating summands do not by
themselves resolve whether the result can vanish.

\subsubsection{An integer recurrence}
\begin{proposition}\label{xx:prop:brec}
For $m\ge2$ and $1\le r\le m$,
\begin{equation}\label{xx:eq:brec}
 b(m,r)=(r-m+1)b(m-1,r)+b(m-1,r-1)+(m-1)b(m-2,r-1).
\end{equation}
Together with rows $b(0,0)=1$ and $(b(1,0),b(1,1))=(0,1)$, this determines
the entire triangle.
\end{proposition}
\begin{proof}
Let $B_r(z)=\Phi(z)^r/r!$. Since
$(1+z)\Phi'(z)=1+z+\Phi(z)$, we have
\[
 (1+z)B_r'(z)=(1+z)B_{r-1}(z)+rB_r(z).
\]
Compare coefficients of $z^{m-1}/(m-1)!$. The result is
\[
 b(m,r)+(m-1)b(m-1,r)
 =b(m-1,r-1)+(m-1)b(m-2,r-1)+r b(m-1,r),
\]
which rearranges to \eqref{xx:eq:brec}.
\end{proof}

For reference, the first rows, omitting column zero, are
\begin{center}
\setlength{\tabcolsep}{4pt}
\small
\begin{tabular}{r|rrrrrrrrr}
$m\backslash r$&1&2&3&4&5&6&7&8&9\\\hline
1&1\\
2&1&1\\
3&$-1$&3&1\\
4&2&$-1$&6&1\\
5&$-6$&0&5&10&1\\
6&24&4&$-15$&25&15&1\\
7&$-120$&$-28$&49&$-35$&70&21&1\\
8&720&188&$-196$&49&0&154&28&1\\
9&$-5040$&$-1368$&944&0&$-231$&252&294&36&1
\end{tabular}
\end{center}

\subsection{The exact zero-count reduction}

\begin{definition}
For $n\ge0$, let
\[
 Z(n)=\#\{(m,r):1\le r\le m\le n,\ b(m,r)=0\}.
\]
Column zero is not included: its zeros for $m>0$ are forced by the definition
and never represent potential nonzero derivative terms.
\end{definition}

\begin{theorem}\label{xx:thm:count}
For every $n\ge0$,
\begin{equation}\label{xx:eq:count}
 \boxed{\displaystyle a(n)=1+\frac{n(n+1)}2-Z(n).}
\end{equation}
\end{theorem}
\begin{proof}
In \eqref{xx:eq:expansion}, the pair $(m,r)$ gives the monomial
\[
 x^{x-(m-r)}(\log x)^{n-m}
\]
with coefficient $\binom nm b(m,r)$. The map
$(m,r)\mapsto(m-r,n-m)$ is injective for fixed $n$: the second coordinate
recovers $m$, and the first then recovers $r$. Thus different pairs do not
cancel each other. The factor $\binom nm$ is positive. There is one nonzero
entry at $(0,0)$, and $m$ potentially nonzero entries $1\le r\le m$ in every
row $m\ge1$. Removing exactly the zero entries proves the formula.
\end{proof}

This argument isolates the main difficulty. Once the triangle is known,
term counting is easy; proving that its list of zeros is complete is not
a consequence of merely writing down the triangle's recurrence.

\subsubsection{Cancellations that are proved}
\begin{proposition}[Symmetry zeros]\label{xx:prop:symmetry}
For every integer $h\ge1$,
\begin{equation}\label{xx:eq:symmetry}
 b(4h+1,2h)=0.
\end{equation}
\end{proposition}
\begin{proof}
Center the product in \eqref{xx:eq:falling} at $r=2h$:
\[
 F_{4h+1}(2h+u)=u\prod_{s=1}^{2h}(u^2-s^2).
\]
This polynomial is odd. Its coefficient of the even power $u^{2h}$ is zero.
\end{proof}

This family is not claimed as new: it appears as Proposition 3.10(iii) in
Jeong \cite{jeong}, with attribution there to Theorem 7 of Lehmer
\cite{lehmer}. Its elementary proof is included to make the reduction
self-contained.

There is one further zero required by the proposed formula:
\begin{equation}\label{xx:eq:exception}
 b(8,5)=0.
\end{equation}
It follows immediately from the factorization proved in
Section~\ref{xx:sec:diagonals}:
\[
 b(m,m-3)=\frac{m(m-1)(m-2)(m-3)(m-5)(m-8)}{48}.
\]
The root $m=5$ is the first symmetry zero; $m=8$ is genuinely additional.

\subsubsection{The precise missing assertion}
\begin{conjecture}[Complete zero classification]\label{xx:conj:zeros}
For integers $1\le r\le m$,
\begin{equation}\label{xx:eq:zeroconj}
 b(m,r)=0\quad\Longleftrightarrow\quad
 (m,r)=(8,5)\ \text{or}\ (m,r)=(4h+1,2h)\text{ for some }h\ge1.
\end{equation}
\end{conjecture}

The right-to-left implication is proved above. The left-to-right implication
is the unproved step in this report.

Let
\[
 \ZZ_0=\{(8,5)\}\cup\{(4h+1,2h):h\ge1\}
\]
be the proved set of zeros, and define
\begin{align*}
 e_m&=\#\{r:1\le r\le m,\ b(m,r)=0,\ (m,r)\notin\ZZ_0\},\\
 E(n)&=\sum_{m=1}^n e_m.
\end{align*}
Thus $E$ counts only zeros not in the proved list.

\begin{theorem}[Exact defect formula]\label{xx:thm:defect}
For every $n\ge0$,
\[
 a(n)=q(n)-E(n).
\]
The following are equivalent: Conjecture~\ref{xx:conj:zeros}; $E(n)=0$ for all
$n$; $a(n)=q(n)$ for all $n$; and the proposed eventual formula
\eqref{xx:eq:eventual} for every $n\ge8$.
\end{theorem}
\begin{proof}
For $n\ge1$, the number of symmetry zeros with first coordinate at most $n$
is $\lfloor(n-1)/4\rfloor$. The zero $(8,5)$ contributes
$\ind{n\ge8}$. Substituting
\[
 Z(n)=\left\lfloor\frac{n-1}{4}\right\rfloor+\ind{n\ge8}+E(n)
\]
into \eqref{xx:eq:count} proves the identity for $n\ge1$; $n=0$ is immediate.
The first three equivalences follow from the definitions. The eventual
formula also forces $E(n)=0$ for every $n\ge8$, and nonnegativity and
monotonicity then force $E(n)=0$ for smaller $n$ as well.
\end{proof}

A hypothetical new zero at row $M$ would decrease every count $a(n)$ for
$n\ge M$ relative to $q(n)$. It would not merely alter a single isolated
term of the sequence. This persistent-deficit property is a useful safeguard
when testing proposed proofs or numerical data.

\subsection{Diagonal polynomials and exact nonvanishing}\label{xx:sec:diagonals}

\subsubsection{A polynomial for each fixed offset}
Define
\begin{equation}\label{xx:eq:H}
 H(z)=\frac{\Phi(z)}z
 =1+\frac z2-\frac{z^2}{6}+\frac{z^3}{12}-\frac{z^4}{20}+\cdots,
 \qquad T_d(Y)=\coef{z}{d}H(z)^Y.
\end{equation}
For a fixed integer $d\ge0$, $T_d$ is a polynomial over $\Q$ in $Y$.

\begin{lemma}\label{xx:lem:diagdegree}
For $m\ge d$,
\begin{equation}\label{xx:eq:diag}
 b(m,m-d)=\fall m d\,T_d(m-d).
\end{equation}
For $d\ge1$, $T_d(Y)$ has degree $d$, leading coefficient $1/(2^d d!)$,
and a factor $Y$. Consequently $T_d(Y)/Y$ is a nonzero polynomial of
degree $d-1$.
\end{lemma}
\begin{proof}
Since $\Phi(z)^{m-d}=z^{m-d}H(z)^{m-d}$, definition \eqref{xx:eq:bdef} gives
\eqref{xx:eq:diag}. Now assume $d\ge1$ and write $h(z)=H(z)-1$. Then
\begin{equation}\label{xx:eq:Tbinomial}
 T_d(Y)=\sum_{\ell=1}^d\binom Y\ell\coef{z}{d}h(z)^\ell.
\end{equation}
Every summand is divisible by $Y$. The only term of degree $d$ in $Y$
comes from $\ell=d$, and $\coef{z}{d}h(z)^d=2^{-d}$. This proves the
remaining assertions.
\end{proof}

The first three cases are
\begin{align}
 b(m,m-1)&=\frac{m(m-1)}2,\label{xx:eq:d1}\\
 b(m,m-2)&=\frac{m(m-1)(m-2)(3m-13)}{24},\label{xx:eq:d2}\\
 b(m,m-3)&=\frac{m(m-1)(m-2)(m-3)(m-5)(m-8)}{48}.
 \label{xx:eq:d3}
\end{align}
For instance,
\[
 T_3(Y)=\frac{Y}{12}-\frac{Y(Y-1)}{12}
                  +\frac{Y(Y-1)(Y-2)}{48}
        =\frac{Y(Y-2)(Y-5)}{48}.
\]
This proves both zeros on the third offset and excludes all others on that
entire diagonal, not just up to a numerical cutoff.

\subsubsection{Prime offsets: an infinite family of nonvanishing results}
For a nonzero rational number $v$, write $v_p(v)$ for its exponent of the
prime $p$ in its prime factorization. A rational number is called
$p$-integral when its denominator in lowest terms is not divisible by $p$.

\begin{theorem}[Prime-offset nonvanishing]\label{xx:thm:prime}
Let $p$ be any prime and let $m\ge p$. Then
\begin{equation}\label{xx:eq:primevaluation}
 b(m,m-p+1)\ne0,\qquad
 v_p\bigl(b(m,m-p+1)\bigr)=v_p\left(\left\lfloor\frac mp\right\rfloor\right).
\end{equation}
In particular, an additional zero cannot occur on any offset $d=m-r$
for which $d+1$ is prime.
\end{theorem}
\begin{proof}
Put $d=p-1$ and $h_j=\coef{z}{j}(H(z)-1)$.
For $1\le j\le p-2$, the number
$h_j=(-1)^{j-1}/(j(j+1))$ is $p$-integral. Moreover,
\[
 p h_{p-1}\equiv1\pmod p.
\]
For odd $p$, this follows from $p h_{p-1}=-1/(p-1)$; for $p=2$, it is
$2h_1=1$.

Divide \eqref{xx:eq:Tbinomial} by $Y$. Its $\ell=1$ summand is $h_{p-1}$.
For every $2\le\ell\le p-1$, a composition of $p-1$ into $\ell$ positive
parts has all parts at most $p-2$. Therefore every coefficient contributing
to $\coef{z}{p-1}h(z)^\ell$ is $p$-integral. Also,
\[
 \frac1Y\binom Y\ell
 =\frac{(Y-1)\cdots(Y-\ell+1)}{\ell!}
\]
has $p$-integral coefficients, because $\ell!$ is not divisible by $p$.
It follows, as a polynomial congruence over the $p$-integral rationals, that
\begin{equation}\label{xx:eq:punit}
 p\,\frac{T_{p-1}(Y)}Y\equiv1\pmod p.
\end{equation}
For every positive integer $Y$, the left side is therefore nonzero and has
$p$-adic valuation zero. Hence
$v_p(T_{p-1}(Y))=v_p(Y)-1$.

Set $Y=m-p+1>0$ and use \eqref{xx:eq:diag}. This yields
\[
 v_p\bigl(b(m,m-p+1)\bigr)=v_p(\fall m p)-1.
\]
Among the $p$ consecutive positive integers $m-p+1,\ldots,m$, exactly one
is divisible by $p$: it is $p\lfloor m/p\rfloor$. Its valuation is
$1+v_p(\lfloor m/p\rfloor)$, proving \eqref{xx:eq:primevaluation}.
\end{proof}

This proof explains, for example, why the residual polynomial $3m-13$ in
\eqref{xx:eq:d2} is a nonzero constant modulo $3$, and why analogous residual
polynomials for offsets $4,6,10,12,16$ are nonzero constants modulo
$5,7,11,13,17$, respectively.

\subsubsection{Finite certificates that cover infinitely many row indices}
For a fixed $d\ge1$, put
\[
 S_d(X)=T_d(X-d).
\]
The factor $X-d$ is forced by Lemma~\ref{xx:lem:diagdegree}. When $d\ge3$ is
odd, the symmetry zero also forces $X-(2d-1)$; for $d=3$ there is in
addition the factor $X-8$.
Remove these factors from $S_d$, clear rational denominators, divide out
the greatest common divisor of the integer coefficients, and normalize the
leading coefficient to be positive. Call the resulting polynomial $R_d(X)$.
Thus, with a positive rational constant $c_d$,
\begin{equation}\label{xx:eq:Rd}
 S_d(X)=c_d(X-d)
 \begin{cases}
 (X-5)(X-8)R_d(X),&d=3,\\
 (X-(2d-1))R_d(X),&d\ge5\text{ odd},\\
 R_d(X),&\text{otherwise}.
 \end{cases}
\end{equation}

\begin{lemma}[Modular certificate]\label{xx:lem:modcert}
If $R\in\Z[X]$ satisfies $R(t)\not\equiv0\pmod p$ for every
$t\in\{0,1,\ldots,p-1\}$, then $R$ has no integer root.
\end{lemma}
\begin{proof}
An integer root, reduced modulo $p$, would be one of those residues and
would give a zero value modulo $p$.
\end{proof}

\begin{theorem}[Computer-assisted diagonal classification]\label{xx:thm:16}
Conjecture~\ref{xx:conj:zeros} holds for every pair $1\le r\le m$ with
$m-r\le16$.
\end{theorem}
\begin{proof}
For offset zero, $b(m,m)=1$. For $1\le d\le16$, compute $S_d$ exactly from
\eqref{xx:eq:Tbinomial} and factor it as in \eqref{xx:eq:Rd}. The supplied
\texttt{data/diagonal\_certificates.json} records the full rational
coefficients of $S_d$, the scalar $c_d$, all removed linear factors, the
full integer coefficients of $R_d$, and every value of $R_d(t)$ modulo
the certifying prime. All of those values are nonzero. The following
summary identifies the certificates:
\begin{center}
\begin{tabular}{rrrr@{\qquad}rrrr}
\toprule
$d$&$\deg R_d$&$p$&roots mod $p$&$d$&$\deg R_d$&$p$&roots mod $p$\\
\midrule
1&0&2&0&9&7&7&0\\
2&1&3&0&10&9&11&0\\
3&0&2&0&11&9&13&0\\
4&3&5&0&12&11&13&0\\
5&3&3&0&13&11&11&0\\
6&5&7&0&14&13&23&0\\
7&5&5&0&15&13&19&0\\
8&7&11&0&16&15&17&0\\
\bottomrule
\end{tabular}
\end{center}
The standard-library program \texttt{code/verify\_diagonals.py} reconstructs
every polynomial using rational arithmetic and checks the factorization
and each residue. The modular lemma excludes every integer root of every
$R_d$. For $m>d$, the removed factor $m-d$ cannot vanish; the remaining
removed factors give exactly the stated symmetry and exceptional zeros.
\end{proof}

These certificates are not extrapolations from finitely many values of
$m$: each modular argument excludes integer roots for \emph{all} $m$ on
its fixed diagonal. What is finite is the set of offsets treated.

For a small example not requiring the machine-readable file,
\begin{align*}
 R_4(X)&=15X^3-390X^2+3245X-8638\equiv2\pmod5,\\
 R_5(X)&=3X^3-103X^2+1118X-3776
                  \equiv2X^2+2X+1\pmod3.
\end{align*}
The second reduced polynomial takes the values $1,2,1$ at the three
residues modulo $3$. Hence there are no additional zeros on offset $4$,
and the only zero on offset $5$ is the symmetry zero $m=9$.

\subsection{The first three columns for all row indices}\label{xx:sec:columns}

The diagonal results keep $m-r$ fixed. A complementary calculation keeps
$r$ fixed and proves nonvanishing in a different infinite region.

\begin{proposition}\label{xx:prop:firstcolumns}
Column $r=1$ has no zero. Column $r=2$ has exactly the zero $m=5$.
Column $r=3$ has no zero.
\end{proposition}
\begin{proof}
For $r=1$, \eqref{xx:eq:falling} gives
\[
 b(1,1)=1,\qquad b(m,1)=(-1)^m(m-2)!\quad(m\ge2).
\]
For $r=2$, put $t=m-3\ge0$. The shifted product is
\[
 F_m(2+u)=u(u+1)(u+2)\prod_{j=1}^{t}(u-j).
\]
The coefficient of $u^2$ is
\begin{equation}\label{xx:eq:col2}
 b(m,2)=(-1)^m(m-3)!\,(2H_{m-3}-3),\qquad m\ge3,
\end{equation}
where $H_t=\sum_{j=1}^t1/j$ and $H_0=0$. The harmonic numbers strictly
increase, and $H_2=3/2$, so the only zero is $m-3=2$. The remaining
boundary entry is $b(2,2)=1$.

For $r=3$, the boundary entry is $b(3,3)=1$. Put $t=m-4\ge0$ and define
\[
 H_t^{(2)}=\sum_{j=1}^t\frac1{j^2},\qquad
 g(t)=6-11H_t+3\bigl(H_t^2-H_t^{(2)}\bigr).
\]
Extracting the coefficient of $u^3$ from
$u(u+1)(u+2)(u+3)\prod_{j=1}^t(u-j)$ gives
\begin{equation}\label{xx:eq:col3}
 b(m,3)=(-1)^t t!\,g(t).
\end{equation}
The elementary identity
\begin{equation}\label{xx:eq:gdiff}
 g(t+1)-g(t)=\frac{6H_t-11}{t+1}
\end{equation}
controls its entire sign pattern. We have $g(0)=6$, $g(1)=-5$, and
$H_3=11/6$. Thus $g$ decreases through $t=3$, has
$g(3)=g(4)=-49/6$, and is strictly increasing thereafter.
Finally, exact rational evaluation gives
\[
 g(19)=-\frac{4571462267}{97772875200}<0,\qquad
 g(20)=\frac{537826687}{1150269120}>0.
\]
It follows that $g(t)<0$ for $1\le t\le19$, whereas $g(t)>0$ for
$t\ge20$. There is no integer zero.
\end{proof}

The last argument illustrates a distinction from a numerical search: exact
values at two adjacent integers, together with a proved monotonicity
statement outside a finite initial segment, settle the entire column.

\subsection{Unconditional quadratic growth}\label{xx:sec:bounds}

\begin{theorem}\label{xx:thm:bounds}
For every integer $n\ge0$,
\begin{equation}\label{xx:eq:bounds}
 n+1+\left\lfloor\frac{n^2}{4}\right\rfloor\le a(n)\le q(n).
\end{equation}
In particular,
\[
 \frac14\le\liminf_{n\to\infty}\frac{a(n)}{n^2}
 \le\limsup_{n\to\infty}\frac{a(n)}{n^2}\le\frac12,
 \qquad a(n)=\Theta(n^2).
\]
\end{theorem}
\begin{proof}
The upper bound is Theorem~\ref{xx:thm:defect}. For the lower bound, count
nonzero entries along diagonals of the triangle in
Theorem~\ref{xx:thm:count}. The main diagonal, including $(0,0)$, contributes
$n+1$ entries.

Fix $1\le d\le n-1$. The diagonal $m-r=d$ has $n-d$ relevant positions,
$m=d+1,\ldots,n$. By Lemma~\ref{xx:lem:diagdegree}, its zero positions are
roots of the nonzero polynomial $T_d(m-d)/(m-d)$ of degree $d-1$.
There are therefore at most $d-1$ such positions, and at least
\[
 \max(0,n-d-(d-1))=\max(0,n-2d+1)
\]
nonzero entries on the diagonal. Summing over all $d\ge1$ gives
\[
 a(n)\ge n+1+\sum_{d\ge1}\max(0,n-2d+1)
       =n+1+\left\lfloor\frac{n^2}{4}\right\rfloor.
\]
The last equality follows by summing the odd numbers $1,3,\ldots,2s-1$
when $n=2s$, and the even numbers $2,4,\ldots,2s$ when $n=2s+1$.
The asymptotic bounds follow immediately.
\end{proof}

\begin{corollary}[A stronger explicit lower bound]\label{xx:cor:primebound}
For every $n\ge0$,
\begin{equation}\label{xx:eq:primebound}
 a(n)\ge n+1+\left\lfloor\frac{n^2}{4}\right\rfloor
       +\sum_{\substack{p\le n\\p\text{ prime}}}\min(p-2,n-p+1).
\end{equation}
\end{corollary}
\begin{proof}
On offset $d=p-1$, Theorem~\ref{xx:thm:prime} guarantees all $n-d=n-p+1$
entries, rather than just $\max(0,n-2p+3)$. The improvement is
$\min(p-2,n-p+1)$. These diagonals are disjoint, so their improvements add.
\end{proof}

The prime-enhanced bound is useful without invoking any theorem about the
density of primes. It still does not close the factor-of-two gap in the
leading quadratic coefficient. In terms of the exact defect,
\begin{equation}\label{xx:eq:asympeq}
 a(n)\sim\frac{n^2}{2}\quad\Longleftrightarrow\quad E(n)=o(n^2).
\end{equation}
Thus the proposed asymptotic equivalent is potentially a weaker target than
the complete zero classification: it would allow a sufficiently sparse set
of additional zeros.

\subsection{The rational function, quasipolynomial, and recurrence}\label{xx:sec:gf}

\subsubsection{Deriving the proposed generating function}
There is one potential term at $(0,0)$ and $m$ potential terms in row $m$.
The cumulative potential count therefore has generating function
\[
 \sum_{n\ge0}\left(1+\frac{n(n+1)}2\right)z^n
 =\frac1{1-z}+\frac{z}{(1-z)^3}.
\]
The symmetry zeros are born at rows $5,9,13,\ldots$, and each remains
absent from every later derivative. Their cumulative generating function is
$z^5/((1-z)(1-z^4))$. The exceptional zero is born at row $8$ and
contributes $z^8/(1-z)$. Subtracting these two \emph{proved} families gives
\begin{align}
 Q(z)
 &=\frac1{1-z}+\frac{z}{(1-z)^3}
       -\frac{z^5}{(1-z)(1-z^4)}-\frac{z^8}{1-z}\notag\\
 &=\frac{1+z^2+z^3+z^6-z^8+z^9+z^{12}-z^{13}}
        {(1-z)^2(1-z^4)}.\label{xx:eq:gfderive}
\end{align}
This explains every correction term in the proposed rational function.

\begin{theorem}[Exact generating-function correction]\label{xx:thm:gfdefect}
With $A(z)=\sum_{n\ge0}a(n)z^n$ and
$\mathcal E(z)=\sum_{m\ge1}e_m z^m$,
\begin{equation}\label{xx:eq:gfdefect}
 A(z)=Q(z)-\frac{\mathcal E(z)}{1-z}.
\end{equation}
Consequently the proposed rational function equals $A(z)$ if and only if
Conjecture~\ref{xx:conj:zeros} is true.
\end{theorem}
\begin{proof}
Sum $a(n)=q(n)-\sum_{m\le n}e_m$ as a formal power series. Each $e_m$
contributes $e_mz^m/(1-z)$ to the cumulative sum.
\end{proof}

\subsubsection{The real form of the complex expression}
For $n\ge8$, division into the four residue classes gives
\begin{equation}\label{xx:eq:residueq}
 q(n)=\frac{n^2}{2}+\frac n4+
 \begin{cases}
 1,&n\equiv0\pmod4,\\
 \tfrac14,&n\equiv1\pmod4,\\
 \tfrac12,&n\equiv2\pmod4,\\
 \tfrac34,&n\equiv3\pmod4.
 \end{cases}
\end{equation}
The same periodic correction is
\[
 \frac{5+(-1)^n+(1-i)(-i)^n+(1+i)i^n}{8},\qquad i^2=-1.
\]
Substitution yields exactly the complex-valued-looking expression printed
in OEIS, whose total value is real and integral. This is an algebraic
identity for $q$, independent of the unresolved equality $a=q$.

For all $n\ge0$, another convenient form is
\[
 q(n)=\frac{n(n+1)}2-\left\lfloor\frac{n-1}{4}\right\rfloor
                         +\ind{1\le n\le7}.
\]
Here $\lfloor-1/4\rfloor=-1$ correctly gives $q(0)=1$.
The first terms are
\[
 1,2,4,7,11,15,21,28,35,43,53,64,76,88,102,117,133,\ldots.
\]

\subsubsection{Why the recurrence starts at fourteen}
Set
\begin{align*}
 D(z)&=(1-z)^2(1-z^4)=1-2z+z^2-z^4+2z^5-z^6,\\
 N(z)&=1+z^2+z^3+z^6-z^8+z^9+z^{12}-z^{13}.
\end{align*}
The identity $D(z)Q(z)=N(z)$ implies, for $n\ge6$,
\begin{equation}\label{xx:eq:inhomogeneous}
 q(n)-2q(n-1)+q(n-2)-q(n-4)+2q(n-5)-q(n-6)
 =\coef{z}{n}N(z).
\end{equation}
The right side vanishes for every $n\ge14$, but at $n=13$ it equals $-1$.
Thus the least starting index from which the homogeneous recurrence holds
for \emph{every} subsequent index of $q$ is $14$.
For $7\le n\le13$, its only nonzero residuals are
\begin{center}
\begin{tabular}{c|rrrr}
$n$&8&9&12&13\\\hline
residual&$-1$&1&1&$-1$
\end{tabular}
\end{center}
This proves the range correction asserted at the start of this part.
If the original derivative sequence satisfies the corrected recurrence
for all $n\ge14$, its directly verified initial values $a(0),\ldots,a(13)$
force $a=q$ by induction. Consequently the corrected recurrence for $a$
is itself equivalent to the main conjecture, once those initial values
are supplied; it is not an independent proved consequence for arbitrary
orders.

The second forward differences $q(n+2)-2q(n+1)+q(n)$ begin
\[
 1,1,1,0,2,1,0,1,\underbrace{2,1,1,0,2,1,1,0,\ldots}_{n\ge8}.
\]
Again, this is proved for $q$; transferring the entire pattern to $a$
requires the missing nonvanishing statement.

\subsection{Reproducible verification}\label{xx:sec:computation}

\subsubsection{What was actually checked}
All reported checks use exact arithmetic. No floating-point threshold was
used to decide that a coefficient was zero.

The C++ implementation with GMP computes the recurrence \eqref{xx:eq:brec}
through row $3000$. It tests all $4{,}501{,}500$ entries with
$1\le r\le m\le3000$. Exactly $750$ are zero: $749$ symmetry zeros and
$(8,5)$. There are no other zeros in that range, and therefore
\begin{equation}\label{xx:eq:3000}
 a(n)=q(n)\quad(0\le n\le3000),\qquad
 a(3000)=4{,}500{,}751.
\end{equation}
These are finite computational results, not an induction proving all $n$.

A separate Python standard-library implementation computes the triangle
through row $1500$ and agrees with the C++ term counts at all $1501$
indices $0,\ldots,1500$. Further independent checks include the complete
polynomials from the differential recurrence \eqref{xx:eq:Prec} through order
$60$, all $496$ shifted-product coefficients \eqref{xx:eq:falling} with
$0\le r\le m\le30$, and $11{,}865$ instances of the prime-adic valuation
formula. All $54$ values displayed directly in the OEIS entry agree with
the exact computation. The separately linked OEIS table through $500$ was
not used as an independent data source for this report.

The sixteen diagonal certificates were reconstructed and checked with a
third, rational-polynomial computation. In addition to showing no unexpected
zeros in a finite triangular region, these certificates prove there are no
unexpected zeros on those fixed diagonals for any row index.

\subsubsection{Algorithm and complexity}
Only two preceding rows are required for \eqref{xx:eq:brec}. A schematic
implementation is
\begin{quote}\ttfamily\small
previousprevious = [1]\\
previous = [0, 1]\\
count = 2\\
for m = 2, ..., N:\\
\hspace*{1em}current[0] = 0\\
\hspace*{1em}for r = 1, ..., m:\\
\hspace*{2em}current[r] = (r-m+1)*previous[r]\\
\hspace*{3em}+ previous[r-1] + (m-1)*previousprevious[r-1]\\
\hspace*{1em}count += number of nonzero current[r] for 1 <= r <= m\\
\hspace*{1em}rotate the three rows
\end{quote}
Indices beyond the end of a row mean zero. The supplied implementations
handle these boundaries explicitly.

The computation uses $O(N^2)$ arithmetic operations and stores $O(N)$
integers. These are not bit-complexity claims: the integers grow with $N$.
From the product formula,
\[
 |b(m,r)|\le\prod_{h=0}^{m-1}(1+|r-h|)\le(m+1)^m,
\]
so a coefficient has $O(m\log m)$ bits. The rolling-row implementation thus
has an $O(N^2\log N)$-bit storage bound. A crude schoolbook arithmetic bound
for the running time is $O(N^3\log^2 N)$ bit operations.

\subsubsection{Files and commands}
The archive contains \texttt{article.tex}, the compiled \texttt{article.pdf},
a build script, both triangle verifiers, the rational certificate verifier,
the complete count and zero-position tables through $3000$, and the exact
verification reports. From the extracted directory, run
\begin{quote}\ttfamily\small
python3 code/verify.py --max-n 1500 --data-dir data\\
python3 code/verify\_diagonals.py --output-dir data --check-only
\end{quote}
The C++ implementation requires a C++17 compiler and GMP development files:
\begin{quote}\ttfamily\small
c++ -O3 -std=c++17 code/verify\_gmp.cpp -lgmpxx -lgmp -o verify\_gmp\\
./verify\_gmp 3000 data
\end{quote}
The Python programs use only the standard library. The diagonal checker
constructs the coefficients afresh rather than trusting the saved residue
tables. Its proof obligations are elementary rational polynomial identities
and finite modular evaluations; it does not require a general-purpose
computer algebra system.

\subsection{What remains to finish the requested proof}

The derivation of the candidate rational function is complete, but the
identification of that function with the derivative count is conditional.
The exact outstanding statement is
\[
 \coef{u}{r}\prod_{h=0}^{m-1}(u+r-h)\ne0
\]
for all $1\le r\le m$ outside the explicitly proved zero set $\ZZ_0$.
The preceding results settle this for $r\le3$, for $m-r\le16$, and whenever
$m-r+1$ is prime, in addition to all rows through $3000$ by exact computation.
They do not cover every remaining pair.

The diagonal-polynomial formulation gives a particularly precise route to
a complete proof. For every offset $d$, it would suffice to show that the
residual polynomial $R_d$ in \eqref{xx:eq:Rd} has no integer root $m>d$.
Irreducibility would be stronger than necessary. A uniform family of
modular nonvanishing certificates would also suffice, but the fact that
such certificates work for $d\le16$ does not establish that they work for
all $d$.

For the asymptotic conjecture alone, an exclusion of every additional zero
is unnecessary: by \eqref{xx:eq:asympeq}, it suffices to bound their cumulative
number by $o(n^2)$. The present root-counting argument proves quadratic
order but not this sharper estimate.

\medskip
\noindent\textbf{Conclusion.}
The printed recurrence with range $n>6$ is disproved, and its correct
range for the proposed sequence is identified. The principal closed formula
is supported by an exact structural reduction, several infinite
nonvanishing theorems, and extensive reproducible verification. A complete
all-order proof of that formula is \emph{not} supplied: the missing converse
in Conjecture~\ref{xx:conj:zeros} is left explicit rather than being replaced
by an inference from numerical agreement.

\clearpage
\section{\texorpdfstring{$x^{x^2}$}{x^(x^2)}: OEIS A290268}\label{xxb:part}

\begin{quote}\small
\textbf{Framework.}  This part instantiates
Theorem~\ref{fw:thm:canonical}\ref{fw:can:b} --- the clause normalized by
$x^{-n}$ --- with the frame $z=x^2$, $L=\log x$;
Proposition~\ref{xxb:prop:recurrence} below is that instance, derived here
independently, and the $-n$ in its recurrence is the normalizer's contribution.
Lemma~\ref{xxb:lem:unique} is the special case of
Theorem~\ref{fw:thm:independence} this part needs.

\smallskip
\textbf{Notation.}  In this part \textbf{$z$ denotes $x^2$}, not $\log x$;
compare Part~\ref{xxc:part}, where $z$ is $\log x$.  See
Section~\ref{fw:notation}.  Throughout this part $a(n)$ denotes the term count
for $x^{x^2}$.

\medskip
\textbf{Abstract of the original report.}\itshape
Let $a(n)$ be the number of nonzero monomials in the fully collected $n$th
derivative of $x^{x^2}$. The conjecture in OEIS A290268 predicts
\[
 a(n)=\U(n):=\frac{n^2}{2}+n+1
 -(n\bmod2)\left(\frac12+\left\lfloor\frac{n+1}{8}\right\rfloor\right).
\]
This article does \upshape\emph{not}\itshape{} establish the equality for all
$n$. It proves the
unconditional bound $a(n)\le\U(n)$ by identifying every cancellation needed to
explain the proposed formula. It also proves positivity throughout the upper
half of the coefficient triangle, classifies all zeros on the first two
logarithmic-deficit diagonals, and gives explicit quadratic lower bounds.
An exact generating-function identity and a finite-difference formula reduce
the full conjecture to a precise nonvanishing assertion. Integer arithmetic
verifies the support through $n=200$; two complete modular computations,
combined with the proved zero identities, certify $a(n)=\U(n)$ for every
$0\le n\le3000$. All source code, finite certificates, and reproduction
instructions accompany the article. The remaining nonvanishing step is
stated explicitly rather than assumed.
\end{quote}

\subsection{The problem and the status of the result}\label{xxb:sec:introduction}
Put
\[
 f(x)=x^{x^2}=\exp(x^2\log x),\qquad x>0.
\]
The sequence A290268 counts the terms in the fully expanded derivative
$f^{(n)}(x)$, after identical monomials have been collected. The OEIS entry
was contributed by Vladimir Reshetnikov in October 2017; its formula section
labels both the generating function and the explicit expression as
conjectural. Peter Luschny's equivalent floor-function form is explicitly
conditional on that conjecture~\cite{oeis290268}.

For clarity, throughout this article the letter $\U$ denotes the
\emph{proposed} answer, not an already proved formula for $a$:
\begin{equation}\label{xxb:eq:U}
 \U(n)=\frac{n^2}{2}+n+1
 -(n\bmod2)\left(\frac12+\left\lfloor\frac{n+1}{8}\right\rfloor\right).
\end{equation}
Equivalently,
\begin{equation}\label{xxb:eq:U-parity}
 \U(2r)=2r^2+2r+1,\qquad
 \U(2r+1)=2(r+1)^2-\left\lfloor\frac{r+1}{4}\right\rfloor.
\end{equation}
The two central issues are different. One must explain why certain
coefficients vanish, and one must also exclude every other cancellation.
The first task is completed here. The second is completed in substantial
regions and in a large finite range, but not in full generality.

\begin{theorem}[Unconditional bounds]\label{xxb:thm:main-bounds}
For every integer $n\ge0$,
\begin{equation}\label{xxb:eq:bounds}
 \Lambda(n)\le a(n)\le\U(n),
 \qquad
 \Lambda(n)=
 \begin{cases}
 (3n^2+10n+8)/8,&n\text{ even},\\[2pt]
 (3n^2+12n+1)/8,&n\text{ odd}.
 \end{cases}
\end{equation}
In particular, $a(n)=\Theta(n^2)$. The proposed leading constant $1/2$ is
not established by these bounds.
\end{theorem}

The upper bound is proved in Section~\ref{xxb:sec:counting}; the lower bound is
proved in Section~\ref{xxb:sec:positivity}. Section~\ref{xxb:sec:computation} gives a
reproducible finite certification of the stronger equality through $3000$.
The exact outstanding assertion appears in Section~\ref{xxb:sec:remaining}.


\subsection{A canonical polynomial expansion}\label{xxb:sec:canonical}
Set $z=x^2$ and $L=\log x$. Define polynomials $P_n\in\Z[z,L]$ by
\begin{equation}\label{xxb:eq:canonical}
 f^{(n)}(x)=f(x)x^{-n}P_n(x^2,\log x),\qquad P_0=1.
\end{equation}
Write
\begin{equation}\label{xxb:eq:c-def}
 P_n(z,L)=\sum_{k,j}c_{n,k,j}z^kL^j.
\end{equation}
Thus $c_{n,k,j}$ is the coefficient of
$x^{x^2+2k-n}(\log x)^j$ in the derivative. Counting terms means counting
nonzero $c_{n,k,j}$, rather than counting the summands generated by a
particular symbolic differentiation algorithm.

\begin{lemma}[Uniqueness of the collected expansion]\label{xxb:lem:unique}
A finite identity $\sum_{k,j}b_{k,j}x^{2k}(\log x)^j=0$ for all $x>0$
forces every $b_{k,j}$ to vanish. Consequently, the expansion
\eqref{xxb:eq:canonical}--\eqref{xxb:eq:c-def} is canonical.
\end{lemma}
\begin{proof}
Substitute $x=e^L$ and group the identity as
$\sum_k e^{2kL}B_k(L)=0$, with polynomial $B_k$. If $K$ is the largest
index with $B_K\ne0$, division by $e^{2KL}$ shows that $B_K(L)$ tends to
zero as $L\to+\infty$: all remaining terms are polynomials multiplied by
decaying exponentials. A nonzero polynomial cannot tend to zero at
$+\infty$. This contradiction eliminates $K$, and then all indices.
\end{proof}

The normalized derivative obeys a particularly small recurrence.
\begin{proposition}[Polynomial and coefficient recurrences]\label{xxb:prop:recurrence}
For every $n\ge0$,
\begin{equation}\label{xxb:eq:P-recurrence}
 P_{n+1}=\bigl(z(2L+1)-n\bigr)P_n
          +2z\frac{\partial P_n}{\partial z}
          +\frac{\partial P_n}{\partial L}.
\end{equation}
With nonexistent coefficients interpreted as zero,
\begin{equation}\label{xxb:eq:c-recurrence}
 \boxed{\;
 c_{n+1,k,j}=(2k-n)c_{n,k,j}+(j+1)c_{n,k,j+1}
             +c_{n,k-1,j}+2c_{n,k-1,j-1}.
 \;}
\end{equation}
For $n\ge1$, the only possible nonzero indices satisfy
\begin{equation}\label{xxb:eq:triangle}
 1\le k\le n,\qquad 0\le j\le k.
\end{equation}
\end{proposition}
\begin{proof}
We have $f'/f=x(2L+1)$, $dz/dx=2z/x$, and $dL/dx=1/x$.
Differentiate \eqref{xxb:eq:canonical} and factor out $f(x)x^{-n-1}$ to obtain
\eqref{xxb:eq:P-recurrence}. Taking the coefficient of $z^kL^j$ gives
\eqref{xxb:eq:c-recurrence}. Starting from $P_0=1$, the first derivative is
$P_1=z(2L+1)$. The four operations in \eqref{xxb:eq:P-recurrence} preserve
$k\ge1$ and $j\le k$, and increase $k$ by at most one. Induction proves
\eqref{xxb:eq:triangle} and integrality.
\end{proof}

For example,
\begin{align*}
 P_1&=z(2L+1),\\
 P_2&=z(2L+3)+z^2(2L+1)^2,\\
 P_3&=2z+z^2(12L^2+24L+9)+z^3(2L+1)^3.
\end{align*}
These have respectively $2$, $5$, and $8$ nonzero coefficients, in
agreement with the definition. Already $P_3$ illustrates a degree loss:
its $z$-coefficient has no $L$-term. Later cancellations are less obvious.
For instance,
\begin{equation}\label{xxb:eq:examples-cancel}
 [z^2]P_7(z,L)=28,\qquad [z^2]P_9(z,L)=-864L.
\end{equation}
The missing $L$-term in the first polynomial and the missing constant in
the second will follow from the same reflection principle.

\subsection{Exact formulas for individual coefficients}\label{xxb:sec:formulas}
\subsubsection{An exponential generating function}
Define two formal power series
\begin{equation}\label{xxb:eq:AB}
 A(t)=t(2+t)=(1+t)^2-1,\qquad
 B(t)=(1+t)^2\log(1+t).
\end{equation}
Formal coefficient extraction is denoted by $[t^n]$. All logarithms in
formal series are expanded at the origin.

\begin{proposition}[Generating function]\label{xxb:prop:egf}
The complete coefficient array is described by
\begin{equation}\label{xxb:eq:egf}
 \sum_{n\ge0}P_n(z,L)\frac{t^n}{n!}
   =\exp\bigl(z(A(t)L+B(t))\bigr).
\end{equation}
In particular, if $m=k-j\ge0$, then
\begin{equation}\label{xxb:eq:coefficient-AB}
 \boxed{\quad
 c_{n,k,j}=\frac{n!}{j!m!}[t^n]A(t)^j B(t)^m.
 \quad}
\end{equation}
\end{proposition}
\begin{proof}
Taylor expansion at the point $x$ gives
\[
 \frac{f(x(1+t))}{f(x)}
 =\sum_{n\ge0}\frac{x^nf^{(n)}(x)}{f(x)}\frac{t^n}{n!}.
\]
The logarithm of the ratio on the left is
\[
 x^2\bigl(((1+t)^2-1)\log x+(1+t)^2\log(1+t)\bigr).
\]
This proves \eqref{xxb:eq:egf} after substitution of $z,L$. Alternatively,
its coefficients satisfy \eqref{xxb:eq:P-recurrence}, so the identity also
holds with $z,L$ algebraically independent. Expanding separately the two
exponentials $\exp(zA(t)L)$ and $\exp(zB(t))$ gives
\eqref{xxb:eq:coefficient-AB}.
\end{proof}

Because both $A$ and $B$ start with a nonzero term of degree one, their
product in \eqref{xxb:eq:coefficient-AB} starts at degree $k$. This independently
explains $k\le n$. The difficulty is that $B$ does not have only positive
coefficients:
\[
 B(t)=t+\frac32t^2+\frac13t^3-\frac1{12}t^4+\frac1{30}t^5-\cdots.
\]
Consequently, the coefficient formula is not itself a proof of
nonvanishing.

\subsubsection{Falling factorials and finite differences}
Let
\[
 F_n(u)=\fall{u}{n}=u(u-1)\cdots(u-n+1),\qquad F_0(u)=1,
\]
and let $\Delta_2G(u)=G(u+2)-G(u)$. Set
$Q_{n,j}(u)=\Delta_2^jF_n(u)$.

\begin{proposition}[Finite-difference formula]\label{xxb:prop:difference}
For $m=k-j\ge0$,
\begin{align}
 c_{n,k,j}
 &=\frac{1}{j!m!}
       \left.\frac{d^m}{du^m}Q_{n,j}(u)\right|_{u=2m}
       \label{xxb:eq:difference}\\
 &=\frac1{j!}[v^m]\sum_{h=0}^j(-1)^{j-h}\binom jh
       F_n(2m+2h+v).\label{xxb:eq:difference-coeff}
\end{align}
\end{proposition}
\begin{proof}
The generalized binomial theorem gives
$n![t^n](1+t)^u=F_n(u)$. Differentiating $m$ times in $u$ introduces
$\log^m(1+t)$. Taking $j$ forward differences of step two introduces
$((1+t)^2-1)^j=A(t)^j$. Evaluation at $u=2m$ introduces
$(1+t)^{2m}$, which combines with the logarithm to give $B(t)^m$.
Equation~\eqref{xxb:eq:coefficient-AB} now proves \eqref{xxb:eq:difference};
Taylor's formula proves \eqref{xxb:eq:difference-coeff}.
\end{proof}

Derivatives of falling factorials also occur in the theory of shifted
Stirling numbers and elementary symmetric harmonic sums; Poon gives a
broader account of these connections~\cite{poon}. No general nonvanishing
theorem from that literature is being assumed here. Formula
\eqref{xxb:eq:difference-coeff} is also an independent way to compute a
coefficient using integer polynomial products.

\subsection{A transformation proving positivity}\label{xxb:sec:positivity}
A fractional-linear change of variable turns a large region of the
coefficient array into a manifestly positive expression. Define
\begin{equation}\label{xxb:eq:H}
 H(u)=\frac{\atanh u}{u}=\sum_{r\ge0}\frac{u^{2r}}{2r+1},
 \qquad H(0)=1.
\end{equation}
The value at zero is understood by its formal power series.

\begin{proposition}[Transformed coefficient formula]\label{xxb:prop:mobius}
For $m=k-j\ge0$ and $n\ge k$,
\begin{equation}\label{xxb:eq:mobius}
 c_{n,k,j}=\frac{n!}{j!m!}\,2^{\,2j+m-n}
 [u^{n-k}](1-u)^{n-2k-1}(1+u)^{2m}H(u)^m.
\end{equation}
Negative integer powers of $1-u$ have their usual formal expansions.
\end{proposition}
\begin{proof}
In the residue expression for $[t^n]A(t)^jB(t)^m$, substitute
$t=2u/(1-u)$. Then
\[
 1+t=\frac{1+u}{1-u},\quad
 A(t)=\frac{4u}{(1-u)^2},\quad
 B(t)=\frac{2(1+u)^2}{(1-u)^2}\atanh u,
 \quad dt=\frac{2\,du}{(1-u)^2}.
\]
Using $[t^n]G(t)=\Res_{t=0}G(t)t^{-n-1}\,dt$, the constant factor becomes
$2^{2j+m-n}$. The remaining power of $1-u$ is $n-2k-1$.
The factors $u^j(\atanh u)^m$ contribute $u^kH(u)^m$, so the resulting
residue extracts the coefficient of $u^{n-k}$. Multiplying by
$n!/(j!m!)$ proves the formula. The formal change-of-variable rule applies
because $2u/(1-u)$ has zero constant term and nonzero linear coefficient.
\end{proof}

\begin{theorem}[A nonvanishing region]\label{xxb:thm:positive}
For $1\le k\le n$ and $0\le j\le k$,
\begin{equation}\label{xxb:eq:positive}
 \begin{split}
 n\le2k&\quad\Longrightarrow\quad c_{n,k,j}>0,\\
 n=2k+1,\ j<k&\quad\Longrightarrow\quad c_{n,k,j}>0.
 \end{split}
\end{equation}
\end{theorem}
\begin{proof}
If $n\le2k$, the factor $(1-u)^{n-2k-1}$ is $(1-u)^{-q}$ with $q\ge1$.
Every coefficient of this factor is strictly positive. The other two
factors in \eqref{xxb:eq:mobius} have nonnegative coefficients and constant
term one. Their product therefore has a strictly positive coefficient in
every nonnegative degree.

If $n=2k+1$, that factor disappears. The condition $j<k$ means $m\ge1$.
The series $H(u)^m$ has a strictly positive coefficient in every even
degree. The polynomial $(1+u)^{2m}$ has positive constant and linear
coefficients. An even target degree can be supplied entirely by $H(u)^m$;
an odd target degree can use the linear term of $(1+u)^{2m}$ and an even
degree from $H(u)^m$. Thus the target coefficient of degree $n-k$ is
strictly positive in both cases. The prefactor in \eqref{xxb:eq:mobius} is
positive.
\end{proof}

\begin{corollary}[Quadratic lower bounds]\label{xxb:cor:lower}
The lower bound in Theorem~\ref{xxb:thm:main-bounds} holds.
\end{corollary}
\begin{proof}
For $n=2r\ge2$, Theorem~\ref{xxb:thm:positive} supplies all coefficients with
$r\le k\le2r$, giving
\[
 a(2r)\ge\sum_{k=r}^{2r}(k+1)=\frac{(r+1)(3r+2)}2.
\]
For $n=2r+1$ with $r\ge1$, it supplies all coefficients with
$r+1\le k\le2r+1$ and the $r$ coefficients $0\le j<r$ at $k=r$.
Hence
\[
 a(2r+1)\ge r+\sum_{k=r+1}^{2r+1}(k+1)
 =\frac{3(r+1)(r+2)}2-1.
\]
These are precisely \eqref{xxb:eq:bounds}. The cases $n=0,1$ follow directly
from $P_0=1$ and $P_1=z(2L+1)$.
\end{proof}

\subsection{The two forced families of zeros}\label{xxb:sec:zeros}
\subsubsection{Degree-forced zeros on the top diagonal}
Setting $j=k$, and therefore $m=0$, in \eqref{xxb:eq:coefficient-AB} yields
an explicit coefficient.
\begin{proposition}[Top logarithmic degree]\label{xxb:prop:top}
For $k\ge1$ and $k\le n\le2k$,
\begin{equation}\label{xxb:eq:top}
 c_{n,k,k}=\frac{n!}{k!}\binom{k}{n-k}2^{2k-n}>0.
\end{equation}
For $n>2k$, one has $c_{n,k,k}=0$.
\end{proposition}
\begin{proof}
The relevant series is the polynomial
$A(t)^k=t^k(2+t)^k$, of degree $2k$. Its coefficient of $t^n$ is
$\binom{k}{n-k}2^{2k-n}$ in the indicated range and zero above it.
\end{proof}
For a fixed $n\ge1$, this removes exactly
$\lfloor(n-1)/2\rfloor$ cells from the possible triangle.

\subsubsection{Reflection-forced cancellations}
The less obvious zeros follow from symmetry about the midpoint of the
roots of a falling factorial.
\begin{lemma}[Reflection of a finite difference]\label{xxb:lem:reflection}
For all $n,j\ge0$,
\begin{equation}\label{xxb:eq:reflection}
 Q_{n,j}(n-2j-1-u)=(-1)^{n-j}Q_{n,j}(u).
\end{equation}
\end{lemma}
\begin{proof}
First,
$F_n(n-1-u)=(-1)^nF_n(u)$, by reversing the order of the factors.
Expand the left side of \eqref{xxb:eq:reflection}:
\[
 \sum_{h=0}^j(-1)^{j-h}\binom jh
 F_n(n-2j-1-u+2h).
\]
Reflection turns its last factor into
$(-1)^nF_n(u+2(j-h))$. Reindexing by $\ell=j-h$ leaves the factor
$(-1)^{n-j}$ times the usual expansion of $\Delta_2^jF_n(u)$.
\end{proof}

\begin{theorem}[Reflection zeros]\label{xxb:thm:reflection-zeros}
Suppose $1\le k\le n$ and $0\le j<k$. Then
\begin{equation}\label{xxb:eq:reflection-zeros}
 k\text{ even},\qquad n=4k-2j+1
 \quad\Longrightarrow\quad c_{n,k,j}=0.
\end{equation}
\end{theorem}
\begin{proof}
Put $m=k-j$. The reflection center in Lemma~\ref{xxb:lem:reflection} is
$(n-2j-1)/2$. Under the displayed condition this equals $2m$, precisely
the evaluation point in \eqref{xxb:eq:difference}. The parity of
$Q_{n,j}(2m+v)$ as a polynomial in $v$ is $n-j$. Its coefficient of $v^m$
therefore vanishes whenever $n-j-m=n-k$ is odd. The relation
$n=4k-2j+1$ makes $n$ odd, so $n-k$ is odd when $k$ is even.
Equation~\eqref{xxb:eq:difference-coeff} completes the proof.
\end{proof}

There is also a direct proof using \eqref{xxb:eq:mobius}. The relation in
\eqref{xxb:eq:reflection-zeros} gives $n-2k-1=2m$, so the series whose
coefficient is extracted becomes
\begin{equation}\label{xxb:eq:even-kernel}
 (1-u^2)^{2m}H(u)^m.
\end{equation}
This series is even. When $k$ is even, $n-k$ is odd, and its coefficient
of that degree is zero. This second proof makes the cancellation visible
without evaluating any large integer coefficient.

\begin{remark}[What symmetry does not say]\label{xxb:rem:symmetry-gap}
The theorem is an implication, not an equivalence. An even power series
has zero odd coefficients, but it can also have zero even coefficients.
Similarly, away from the reflection center, a derivative of a polynomial
can vanish for reasons unrelated to symmetry. Excluding such additional
zeros is a separate requirement of the full conjecture.
\end{remark}

\subsection{Counting the forced zeros}\label{xxb:sec:counting}
For $n\ge1$, the triangle \eqref{xxb:eq:triangle} has
\begin{equation}\label{xxb:eq:triangle-size}
 T(n)=\sum_{k=1}^n(k+1)=\frac{n(n+3)}2
\end{equation}
cells. Let $Z_n^{\mathrm{top}}$ be the top-diagonal cells removed by
Proposition~\ref{xxb:prop:top}, and let $Z_n^{\mathrm{ref}}$ be the cells in
Theorem~\ref{xxb:thm:reflection-zeros}. These sets are disjoint, since the
second requires $j<k$.

\begin{lemma}[Number of reflection zeros]\label{xxb:lem:reflection-count}
For even $n$, $Z_n^{\mathrm{ref}}$ is empty. For odd $n$,
\begin{equation}\label{xxb:eq:reflection-count}
 |Z_n^{\mathrm{ref}}|=\left\lfloor\frac{n+1}{8}\right\rfloor.
\end{equation}
\end{lemma}
\begin{proof}
Reflection zeros have odd $n$. Write $n=2r+1$ and $k=2\ell$.
The reflection relation determines
$j=2k-r=4\ell-r$. The inequalities $0\le j<k$ become
\[
 \frac r4\le\ell<\frac r2.
\]
For $r\ge1$, these automatically give $\ell\ge1$ whenever the interval
contains an integer. Its number of integers is
\[
 \left\lceil\frac r2\right\rceil-
 \left\lceil\frac r4\right\rceil
 =\left\lfloor\frac{r+1}{4}\right\rfloor,
\]
as is checked in the four residue classes of $r$ modulo four. For $r=0$
both sides are zero. Since $(r+1)/4=(n+1)/8$, the result follows.
\end{proof}

\begin{theorem}[The conjectured expression is an upper bound]\label{xxb:thm:upper}
For every $n\ge0$,
\begin{equation}\label{xxb:eq:upper}
 a(n)\le\U(n).
\end{equation}
Equality holds for a given $n$ if and only if the coefficient triangle has
no zeros other than $Z_n^{\mathrm{top}}\cup Z_n^{\mathrm{ref}}$.
\end{theorem}
\begin{proof}
For even $n=2r\ge2$, subtraction of the top-diagonal zeros from
\eqref{xxb:eq:triangle-size} gives
\[
 2r^2+3r-(r-1)=2r^2+2r+1.
\]
For odd $n=2r+1$, subtraction of both disjoint families gives
\[
 2r^2+5r+2-r-\left\lfloor\frac{r+1}{4}\right\rfloor
 =2(r+1)^2-\left\lfloor\frac{r+1}{4}\right\rfloor.
\]
These are \eqref{xxb:eq:U-parity}. Additional zeros could only decrease the
number of nonzero coefficients, proving the inequality and the stated
equivalence. Finally $a(0)=\U(0)=1$.
\end{proof}

For example, at $n=7$ there are $35$ potential cells, three top-diagonal
zeros, and the reflection zero $(k,j)=(2,1)$, leaving $31$. At $n=9$,
there are $54$ potential cells, four top-diagonal zeros, and the reflection
zero $(2,0)$, leaving $49$. The exact expansions in
\eqref{xxb:eq:examples-cancel} display those two different cancellations.
The floor term in the proposed formula is thus explained by a proved
symmetry, not merely by fitting initial values.

\subsection{Nonvanishing on the first two deficit diagonals}\label{xxb:sec:small-m}
Call $m=k-j$ the \emph{logarithmic deficit}. This section gives a complete
zero classification for $m=1$ and $m=2$. Together with
Theorem~\ref{xxb:thm:positive}, it narrows the unproved portion of the problem.

\subsubsection{A beta-integral representation}
Suppose $n>2k$, $m=k-j\ge1$, and put
\begin{equation}\label{xxb:eq:d}
 d=n-2k-1\ge0.
\end{equation}
For sufficiently small complex $\varepsilon$, define
\begin{equation}\label{xxb:eq:I}
 I_{m,j,d}(\varepsilon)=
 \int_0^1 t^{2m+\varepsilon}(1-t)^{d-\varepsilon}(2t-1)^j\,dt.
\end{equation}
The integral and all derivatives required below converge uniformly on
compact subsets of $|\varepsilon|<1/2$.

\begin{lemma}[Integral formula]\label{xxb:lem:beta}
Under these assumptions,
\begin{equation}\label{xxb:eq:beta-Q}
 Q_{n,j}(2m+\varepsilon)
 =(-1)^{n-1}n!\frac{\sin(\pi\varepsilon)}{\pi}
    I_{m,j,d}(\varepsilon),
\end{equation}
and consequently
\begin{equation}\label{xxb:eq:beta-c}
 c_{n,k,j}=\frac{(-1)^{n-1}n!}{j!}
 [\varepsilon^m]\left(
 \frac{\sin(\pi\varepsilon)}{\pi}I_{m,j,d}(\varepsilon)\right).
\end{equation}
\end{lemma}
\begin{proof}
Euler's reflection formula and beta integral~\cite{dlmf-reflection,dlmf-beta}
give, initially away from integer $\alpha$,
\begin{align*}
 F_n(\alpha)
 &=(-1)^{n-1}\frac{\sin(\pi\alpha)}{\pi}
      \Gamma(\alpha+1)\Gamma(n-\alpha)\\
 &=(-1)^{n-1}n!\frac{\sin(\pi\alpha)}{\pi}
      \int_0^1t^\alpha(1-t)^{n-\alpha-1}\,dt.
\end{align*}
In the $j$th difference put $\alpha=2m+2h+\varepsilon$. The sine factor
is independent of the integer $h$. The finite sum inside the integral is
\[
 \sum_{h=0}^j(-1)^{j-h}\binom jh
     \left(\frac{t}{1-t}\right)^{2h}
 =\frac{(2t-1)^j}{(1-t)^{2j}}.
\]
The remaining exponent of $1-t$ is $d-\varepsilon$, proving
\eqref{xxb:eq:beta-Q}. All terms are analytic near $\varepsilon=0$, so the
identity extends to zero as well. Taking the coefficient of
$\varepsilon^m$ and using \eqref{xxb:eq:difference-coeff} gives
\eqref{xxb:eq:beta-c}.
\end{proof}

\subsubsection{Deficit one}
\begin{theorem}[Complete signs for $m=1$]\label{xxb:thm:m1}
Suppose $m=k-j=1$ and $n>2k$, and define $d$ by \eqref{xxb:eq:d}. Then
\begin{equation}\label{xxb:eq:m1-sign}
 \sgn c_{n,k,j}=(-1)^{n-1}
 \begin{cases}
 1,&j\text{ even},\\
 \sgn(2-d),&j\text{ odd}.
 \end{cases}
\end{equation}
In particular the only zeros in this range have $j$ odd and $d=2$.
\end{theorem}
\begin{proof}
The coefficient of $\varepsilon$ in \eqref{xxb:eq:beta-c} is
\begin{equation}\label{xxb:eq:m1-integral}
 c_{n,k,j}=\frac{(-1)^{n-1}n!}{j!}
 \int_0^1t^2(1-t)^d(2t-1)^j\,dt.
\end{equation}
For even $j$, the integrand is nonnegative and positive except at finitely
many points, so the integral is strictly positive. For odd $j$, pair $t$
with $1-t$ and integrate over $1/2<t<1$. The paired integrand is
\[
 (2t-1)^j\bigl(t^2(1-t)^d-t^d(1-t)^2\bigr).
\]
The ratio of the two positive terms in parentheses is
$(t/(1-t))^{2-d}$. Since $t/(1-t)>1$, their difference is identically
zero when $d=2$ and otherwise has the strict sign of $2-d$ throughout
the open interval. This proves the sign formula.
\end{proof}
These zeros are exactly the previously identified reflection zeros:
$j$ odd means $k=j+1$ is even, and $d=2$ means
$n=2k+3=4k-2j+1$.

\subsubsection{Deficit two}
\begin{theorem}[Complete signs for $m=2$]\label{xxb:thm:m2}
Suppose $m=k-j=2$ and $n>2k$. Then
\begin{equation}\label{xxb:eq:m2-sign}
 \sgn c_{n,k,j}=(-1)^{n-1}
 \begin{cases}
 1,&j\text{ odd},\\
 \sgn(4-d),&j\text{ even}.
 \end{cases}
\end{equation}
The only zeros in this range have $j$ even and $d=4$.
\end{theorem}
\begin{proof}
The coefficient of $\varepsilon^2$ in \eqref{xxb:eq:beta-c} is
\begin{equation}\label{xxb:eq:m2-integral}
 c_{n,k,j}=\frac{(-1)^{n-1}n!}{j!}
 \int_0^1t^4(1-t)^d(2t-1)^j
       \log\frac{t}{1-t}\,dt.
\end{equation}
For odd $j$, the factors $(2t-1)^j$ and $\log(t/(1-t))$ have the same
sign, so the integral is strictly positive. For even $j$, their product
is antisymmetric under $t\mapsto1-t$. Pairing the integral as before
leaves, on $1/2<t<1$,
\[
 (2t-1)^j\log\frac{t}{1-t}
 \bigl(t^4(1-t)^d-t^d(1-t)^4\bigr).
\]
The first two factors are positive there. The bracket has the strict
sign of $4-d$, or is identically zero when $d=4$. This proves the claim.
\end{proof}
Again the zero condition is exactly reflection: $k=j+2$ is even and
$n=2k+5=4k-2j+1$.

For $n\le2k$, positivity has already been proved, and the borderline
$n=2k+1$ is included in both sign arguments. Thus these two theorems,
together with Proposition~\ref{xxb:prop:top}, classify all zeros for
$k-j\le2$ throughout the possible triangle.

\subsubsection{Why the same sign proof does not settle higher deficits}
At $m=3$, formula \eqref{xxb:eq:beta-c} gives the integral of
\begin{equation}\label{xxb:eq:m3-kernel}
 t^6(1-t)^d(2t-1)^j
 \left(\frac12\log^2\frac{t}{1-t}-\frac{\pi^2}{6}\right)
\end{equation}
multiplied by $(-1)^{n-1}n!/j!$. The last factor changes sign inside
$(0,1)$. Thus positivity of the weight, or pairing $t$ with $1-t$, no
longer excludes cancellation. This is a concrete obstruction to simply
extending the proofs above. It is not a counterexample to the conjecture.

\subsection{Generating functions and equivalent formulas for the upper bound}
\label{xxb:sec:gf}
All identities in this section are proved identities for $\U(n)$.
They become identities for $a(n)$ only after the remaining nonvanishing
assertion is supplied, or when restricted to a certified finite range.

\subsubsection{Derivation of the rational generating function}
By \eqref{xxb:eq:U-parity}, the even and odd subsequences have generating
functions
\begin{align}
 E(y)&=\sum_{r\ge0}\U(2r)y^r
       =\frac{(1+y)^2}{(1-y)^3},\label{xxb:eq:E}\\
 O(y)&=\sum_{r\ge0}\U(2r+1)y^r
       =\frac{2(1+y)}{(1-y)^3}
          -\frac{y^3}{(1-y)(1-y^4)}.\label{xxb:eq:O}
\end{align}
For the last term, write
$\lfloor(r+1)/4\rfloor=\sum_{h\ge1}\mathbf1_{r\ge4h-1}$ and sum each
geometric tail. Therefore
\begin{align}
 G(t):=\sum_{n\ge0}\U(n)t^n
 &=E(t^2)+tO(t^2)\notag\\
 &=\frac{1+t^2}{(1-t)^3(1+t)}
       -\frac{t^7}{(1-t^2)(1-t^8)}\notag\\
 &=\frac{N(t)}{(1-t)(1-t^2)(1-t^8)},\label{xxb:eq:GF}
\end{align}
where
\begin{equation}\label{xxb:eq:N}
 N(t)=1+t+2t^2+2t^3+2t^4+2t^5+2t^6+t^7+2t^8+t^9.
\end{equation}
Indeed,
$N(t)=(1+t^2)(1+t+\cdots+t^7)-t^7(1-t)$.
This is exactly the rational function proposed in the OEIS entry~\cite{oeis290268}.
The derivation shows why period eight appears: it counts the even values
of $k$ on the reflection line, in addition to the parity split of the
ordinary triangle.

\subsubsection{The eight constituent polynomials}
Writing $n=8q+s$ with $0\le s\le7$, \eqref{xxb:eq:U} becomes the following
degree-two quasipolynomial.
\begin{center}
\begin{tabular}{cl}
\toprule
$s$ & $\U(8q+s)$\\
\midrule
0 & $32q^2+8q+1$\\
1 & $32q^2+15q+2$\\
2 & $32q^2+24q+5$\\
3 & $32q^2+31q+8$\\
4 & $32q^2+40q+13$\\
5 & $32q^2+47q+18$\\
6 & $32q^2+56q+25$\\
7 & $32q^2+63q+31$\\
\bottomrule
\end{tabular}
\end{center}
The trigonometric expression from the OEIS entry is another way to write
the same eight polynomials:
\begin{equation}\label{xxb:eq:trig}
 \U(n)=\frac{8n^2+15n+14+(n+2)(-1)^n
 +(2-4\sqrt2\sin(\pi n/4))\sin(\pi n/2)}{16}.
\end{equation}
To verify the equivalence without an approximation, substitute $n=8q+s$.
The sine factors take their exact values at multiples of $\pi/4$ and
$\pi/2$. Substitution in each of the eight residue classes yields the
polynomials in the table. In particular, no irrational value remains.

\subsubsection{Linear recurrences}
Multiplication of \eqref{xxb:eq:GF} by its denominator yields
\begin{equation}\label{xxb:eq:rec11}
 \begin{split}
 \U(n)={}&\U(n-1)+\U(n-2)-\U(n-3)+\U(n-8)\\
         &-\U(n-9)-\U(n-10)+\U(n-11),\qquad n\ge11.
 \end{split}
\end{equation}
Multiplying the denominator by $1+t$ gives the convenient order-twelve
form recorded in the OEIS entry:
\begin{equation}\label{xxb:eq:rec12}
 \U(n)=2\U(n-2)-\U(n-4)+\U(n-8)-2\U(n-10)+\U(n-12),
 \qquad n\ge12.
\end{equation}
The twelve initial values are
\[
 1,2,5,8,13,18,25,31,41,49,61,71.
\]
Using these recurrences to \emph{generate} the claimed values of $a(n)$
would be circular in a proof of the conjecture. The computations in the
next section instead use the independently derived coefficient recurrence
\eqref{xxb:eq:c-recurrence}.

\subsection{Exact finite verification and certificates}\label{xxb:sec:computation}
\subsubsection{Integer arithmetic and independent formula checks}
The accompanying Python program applies \eqref{xxb:eq:c-recurrence} with
arbitrary-precision integers. Through $n=200$, it checks not only the
term count but the entire zero pattern: a coefficient is zero if and only
if it belongs to one of the two proved zero families. This examines
$1{,}373{,}500$ cells with $1\le k\le n\le200$, $0\le j\le k$.

For $0\le n\le16$, the program independently recomputes $969$ triples
$(n,k,j)$, including the $k=0$ boundary, using each of
\eqref{xxb:eq:coefficient-AB}, \eqref{xxb:eq:difference-coeff}, and
\eqref{xxb:eq:mobius}. Fractions are exact rational numbers. All three formulas
agree with the integer recurrence. The positivity and deficit-one and
-deficit-two sign theorems are also checked throughout the integer range.
These checks test the derivations and their implementation; the proofs in
the preceding sections do not depend on them.

\subsubsection{Why modular arithmetic gives a rigorous finite certificate}
A nonzero residue modulo any integer $M\ge2$ proves that the original
integer is nonzero. A zero residue, in contrast, might mean only that
$M$ divides a nonzero integer. We use two moduli,
\begin{equation}\label{xxb:eq:moduli}
 M_1=1{,}000{,}000{,}007,\qquad M_2=1{,}000{,}000{,}009.
\end{equation}
Their primality is not needed for this argument: only exact reduction of
an integer recurrence is used, with no modular division.

For each modulus, a C++ program computes every row through $n=3000$ and
records every zero residue outside the two mathematically proved zero
families. There are three such residues modulo $M_1$ and four modulo
$M_2$. The two lists are disjoint. Their union and the corresponding
residues are as follows.
\begin{center}
\small
\begin{tabular}{rrr|rr}
\toprule
$n$ & $k$ & $j$ & $c_{n,k,j}\bmod M_1$ & $c_{n,k,j}\bmod M_2$\\
\midrule
1471&1046&521&847578491&0\\
1701&610&505&0&439228697\\
1854&1713&549&0&366891231\\
2336&1689&145&0&782191592\\
2543&1785&663&405694396&0\\
2807&2566&502&607368440&0\\
2842&2673&2118&474540977&0\\
\bottomrule
\end{tabular}
\end{center}
Every entry in this table is therefore a nonzero integer coefficient.
All other cells outside the proved zero families have a nonzero residue
in both runs. Each modulus run checks $4{,}509{,}002{,}501$ potential
cells, including the single cell at $n=0$.

\begin{proposition}[Finite computational certification]\label{xxb:prop:finite}
The accompanying exact computations, together with
Proposition~\ref{xxb:prop:top} and Theorem~\ref{xxb:thm:reflection-zeros}, certify
\begin{equation}\label{xxb:eq:finite}
 a(n)=\U(n)\qquad\text{for every integer }0\le n\le3000.
\end{equation}
\end{proposition}
\begin{proof}[Certificate argument]
The two analytic results prove that every predicted zero is an actual
integer zero. For any other coefficient in the finite range, membership
in neither, one, or both modular zero lists exhausts the possibilities.
The intersection of the lists is empty. Thus at least one of its two
residues is nonzero, proving that the integer coefficient is nonzero.
Theorem~\ref{xxb:thm:upper} now gives the equality for every row in the range.
\end{proof}
This is a deterministic finite verification, not a probabilistic claim
based on the size of the moduli. It is also not a formal proof-assistant
certificate: it relies on the supplied ordinary programs and their
execution, which can be independently reproduced.

\subsubsection{Implementation details and limitations}
The C++ program stores two coefficient rows in arrays of unsigned
32-bit integers. An update is evaluated in signed 64-bit arithmetic
before reduction. Its explicit input bounds guarantee that these
intermediates cannot overflow. It never approximates logarithms or
transcendental values. The arrays have $O(N^2)$ entries and the computation
uses $O(N^3)$ fixed-modulus arithmetic operations through order $N$.

The Python implementation uses a different update organization: it
scatters each existing coefficient into its four successors, whereas
the C++ implementation gathers four contributions into each destination.
The small-range generating-function and finite-difference calculations
provide further independent comparisons. The archive retains both raw
modular zero lists, row-by-row support statistics, explicit nonzero
witnesses, the combined certificate, and the exact small-range counts.

No finite range, regardless of its size, proves the nonvanishing assertion
for all $n$. In particular, the certificate supplies neither an induction
step beyond $3000$ nor a theorem bounding where a first additional zero
could occur.

\subsection{The precise missing step}\label{xxb:sec:remaining}
The coefficient problem can now be stated independently of the original
transcendental function.
\begin{conjecture}[Complete support classification]\label{xxb:conj:support}
Let $n\ge1$, $1\le k\le n$, $0\le j\le k$, and $m=k-j$. Then
\begin{equation}\label{xxb:eq:classification}
 \frac{1}{j!m!}
 \left.\frac{d^m}{du^m}\Delta_2^j\fall{u}{n}\right|_{u=2m}=0
\end{equation}
if and only if at least one of the following holds:
\begin{equation}\label{xxb:eq:only-zeros}
 \bigl(j=k\text{ and }n>2k\bigr),
 \qquad
 \bigl(k\text{ even and }n=4k-2j+1\bigr).
\end{equation}
\end{conjecture}
The implication from \eqref{xxb:eq:only-zeros} to \eqref{xxb:eq:classification}
is proved above. The converse remains unproved in the following precise
region:
\begin{equation}\label{xxb:eq:gap-region}
 m=k-j\ge3,\qquad n\ge2k+2,
\end{equation}
away from the reflection zeros. All other regions have been handled by
the preceding analytic theorems.

\begin{theorem}[Exact reduction of the OEIS conjecture]\label{xxb:thm:equivalence}
The equality $a(n)=\U(n)$ for every $n\ge0$ is equivalent to
Conjecture~\ref{xxb:conj:support}. It is also equivalent to proving the
nonvanishing assertion only in the remaining region
\eqref{xxb:eq:gap-region}, excluding the reflection-zero family.
\end{theorem}
\begin{proof}
Formula~\eqref{xxb:eq:difference} identifies the expression in
\eqref{xxb:eq:classification} with the actual derivative coefficient.
Theorem~\ref{xxb:thm:upper} counts exactly the cells left after removal of the
two proved, disjoint zero families. Hence any additional zero makes the
corresponding inequality strict, and no additional zeros makes it an
equality. Theorem~\ref{xxb:thm:positive}, Proposition~\ref{xxb:prop:top}, and
Theorems~\ref{xxb:thm:m1}--\ref{xxb:thm:m2} eliminate the complementary regions.
\end{proof}

A tempting shortcut would assert that a derivative of a falling factorial
cannot vanish at an integer other than its reflection center. That
assertion is false. Direct differentiation gives
\begin{equation}\label{xxb:eq:false-shortcut}
 \frac{d^5}{du^5}\fall{u}{8}
 =3360(u-5)(u-2)(2u-7).
\end{equation}
It vanishes at $u=2$ and $u=5$, as well as at the reflection center
$u=7/2$. This is \emph{not} a counterexample to
Conjecture~\ref{xxb:conj:support}, because the derivative order, evaluation
point, and finite-difference order in that conjecture are linked in a
special way. It does show why those constraints must be used in a genuine
proof of the missing implication.

The transformed formula provides an alternative target. With
$d=n-2k-1\ge1$, the outstanding task is to prove that
\begin{equation}\label{xxb:eq:gap-transform}
 [u^{k+d+1}](1-u)^d(1+u)^{2m}
       \left(\frac{\atanh u}{u}\right)^m
\end{equation}
is nonzero for $m\ge3$, except when $d=2m$ and $k$ is even. The
reflection case is transparent because the series then becomes even;
the remaining finite-difference cancellations are the essential issue.
Equation~\eqref{xxb:eq:gap-transform} is an equivalent residual lemma, not an
additional theorem claimed by this article.

\subsection{Conclusion}
The proposed formula has a rigorous structural explanation. The obvious
degree loss removes the top logarithmic coefficients with $n>2k$.
A separate reflection symmetry removes exactly
$\lfloor(n+1)/8\rfloor$ further coefficients for odd $n$. Their total
produces the OEIS floor expression and its rational generating function
as an unconditional upper bound.

A fractional-linear transformation proves that a large region of the
coefficient triangle is strictly positive, and beta-integral arguments
exclude additional zeros on the first two deficit diagonals. These
results prove quadratic growth and reduce the full formula to the
explicit residual nonvanishing problem \eqref{xxb:eq:gap-transform}.
Exact finite computations settle every derivative order through $3000$.
The global equality is not proved here: its remaining obligation is the
nonvanishing statement in \eqref{xxb:eq:gap-region}, not any of the counting
or generating-function manipulations.

\subsection{Reproducing the computations}\label{xxb:app:reproduce}
The archive needs Python 3.10 or later and a C++17 compiler for the
computations. Python uses only its standard library. The PDF can be
rebuilt with a conventional LaTeX installation containing the packages
listed in the source preamble.

Run the following commands from the archive directory. The Makefile
wraps the same operations.
\begin{lstlisting}
python3 code/verify_exact.py --max-n 200 --formula-n 16
mkdir -p build
g++ -O3 -std=c++17 -Wall -Wextra \
    code/verify_modular.cpp -o build/verify_modular
build/verify_modular 3000 1000000007 data/p1000000007 \
    data/watch_coordinates.csv
build/verify_modular 3000 1000000009 data/p1000000009 \
    data/watch_coordinates.csv
python3 code/combine_certificates.py
pdflatex -interaction=nonstopmode -halt-on-error \
    A290268_partial_results.tex
pdflatex -interaction=nonstopmode -halt-on-error \
    A290268_partial_results.tex
\end{lstlisting}
The watch file requests explicit residues at the seven exceptional
coordinates found in the modular runs. It does not restrict the search:
the C++ program checks every coefficient in every row regardless of the
watch file. The combiner checks that the complete exceptional zero lists
are disjoint and that all listed witnesses are consistent with them.

For a smaller test, the optional watch file can be omitted. For example,
\texttt{build/verify\_modular 100 1000000007 data/test} performs the full
modular recurrence through $100$. A watch file containing indices beyond
the chosen bound is rejected, rather than silently ignored.

\subsection{Files included in the archive}
\begin{longtable}{@{}>{\raggedright\arraybackslash}p{0.43\textwidth}>{\raggedright\arraybackslash}p{0.53\textwidth}@{}}
\toprule
File or directory & Purpose\\
\midrule
\endhead
\texttt{article.tex} & Complete source of the merged article; this part is its Part~3.\\
\texttt{article.pdf} & The compiled merged article.\\
\texttt{README.md}, \texttt{Makefile} & Status, dependency notes, and reproduction commands.\\
\texttt{code/verify\_exact.py} & Arbitrary-precision recurrence and three independent coefficient-formula checks.\\
\texttt{code/verify\_modular.cpp} & Complete modular coefficient search with explicit overflow bounds.\\
\texttt{code/combine\_certificates.py} & Verification and combination of the two finite modular certificates.\\
\texttt{data/exact\_counts.csv} & Exact term counts through $200$.\\
\texttt{data/certified\_counts.csv} & Certified term counts through $3000$.\\
\texttt{data/p1000000007\_*} and \texttt{data/p1000000009\_*} & Complete exceptional zero lists, row statistics, witness residues, run logs, and summaries.\\
\texttt{data/modular\_witnesses.csv} & The seven coordinates at which one residue vanishes, with the nonzero residue in the other modulus.\\
\texttt{data/*summary.json} & Machine-readable verification scope and results.\\
\texttt{data/sample\_coefficients.json} & Fully collected low-order coefficient arrays and cancellation examples.\\
\bottomrule
\end{longtable}

\clearpage
\section{\texorpdfstring{$x^{x^x}$}{x^(x^x)}: OEIS A281434}\label{xxc:part}

\begin{quote}\small
\textbf{Framework.}  This part instantiates
Theorem~\ref{fw:thm:canonical}\ref{fw:can:a} with the three-variable frame
$u=x^x$, $v=x^{-1}$, $z=\log x$; Proposition~\ref{xxc:prop:recurrence} below is
that instance, derived here independently, and it is the only one of the three
whose frame is not two-dimensional.  Because $u=x^x$ is not of the form
$x^{\lambda}$ with $\lambda$ constant, the independence statement this part
needs is \emph{not} a special case of Theorem~\ref{fw:thm:independence}; the
asymptotic argument that replaces it is given in full below.

\smallskip
\textbf{Notation.}  In this part \textbf{$z$ denotes $\log x$}, not $x^2$;
compare Part~\ref{xxb:part}.  See Section~\ref{fw:notation}.  Throughout this
part $a(n)$ denotes the term count for $x^{x^x}$.

\medskip
\textbf{Abstract of the original report.}\itshape
Let $a(n)$ be the number of nonzero monomials in the fully expanded $n$th
derivative of $x^{x^x}$, after collecting identical terms. We replace repeated
symbolic differentiation by a seven-transition integer recurrence on three
exponents. It computes the prefix through $n$ in $O(n^4)$ integer arithmetic
operations and $O(n^3)$ live integers. A faster implementation performs the
recurrence modulo a machine-size integer and verifies every possible zero
by an independent exact coefficient formula; this method is deterministic,
not a probabilistic zero test. We prove the explicit bounds
\[
 \frac{7n^3+39n^2+26n+3\epsn(n-1)}{24}
 \ \le a(n)\le\ \frac{2n^3+3n^2+4n}{3},\qquad n\ge1,
\]
where $\epsn$ is $1$ for odd $n$ and $0$ for even $n$. In particular,
$a(n)=\Theta(n^3)$. The upper bound counts an explicit lattice envelope;
the lower bound uses nonvanishing leading coefficients and high-multiplicity
roots of coefficient polynomials. We also identify an unavoidable family of
logarithm-free cancellations and an infinite family with positive logarithm
exponent. The stronger statement $a(n)\sim \tfrac23n^3$, and the still stronger
linear-deficit conjecture, are \upshape\emph{not proved here}\itshape. Exact
code, tests, terms through $n=100$, and selected computations through $n=200$
accompany this article.
\end{quote}

\subsection{The problem and the scope of the results}

The sequence A281434 is defined in the OEIS as the number of terms in the fully
expanded $n$th derivative of $x^{x^x}$, with offset $0$~\cite{oeis281434}.
Exponentiation is right-associated. Thus the function in question is
\[
 f(x)=x^{(x^x)}=\exp\bigl(x^x\log x\bigr),\qquad x>0,
\]
not $(x^x)^x=x^{x^2}$. Its initial term counts are
\[
1,3,12,30,64,113,188,285,415,577,780,1017,1312,\ldots.
\]
The positivity restriction supplies an unambiguous real logarithm; all
subsequent coefficient manipulations are exact polynomial identities.

The phrase ``fully expanded'' must include collecting like terms and deleting
zero coefficients. Counting product-rule branches is a different problem.
Even recording every monomial that could be reached by differentiation gives
an overcount: distinct contributions can cancel exactly.

For $n\ge1$, define
\begin{equation}\label{xxc:eq:B}
 B(n)=\frac{2n^3+3n^2+4n}{3},\qquad
 \Delta(n)=B(n)-a(n).
\end{equation}
The polynomial $B(n)$ is A037237$(n-1)$~\cite{oeis037237}. The related OEIS
entry A352697 records the differences $\Delta(n)$ and conjectures both
$\Delta(n)=O(n)$ and infinitely many zero differences~\cite{oeis352697}.
These observations are useful motivation, but neither is needed by our
algorithms or proofs.

\begin{theorem}[Main results]\label{xxc:thm:main}
There is a deterministic exact algorithm for A281434 requiring $O(n^4)$
integer arithmetic operations and $O(n^3)$ stored integers to compute
$a(0),\ldots,a(n)$. There is also a machine-modular implementation for a
single $a(n)$, with exact verification of every candidate cancellation.
For $n\ge1$,
\begin{equation}\label{xxc:eq:mainbounds}
 L(n):=\frac{7n^3+39n^2+26n+3\epsn(n-1)}{24}
 \le a(n)\le B(n),
 \quad \epsn=\frac{1-(-1)^n}{2}.
\end{equation}
Consequently
\begin{equation}\label{xxc:eq:theta}
 a(n)=\Theta(n^3),\qquad
 \lim_{n\to\infty}\frac{\log a(n)}{\log n}=3.
\end{equation}
\end{theorem}

The distinction between a growth \emph{order} and an asymptotic
\emph{equivalent} is important. Equation~\eqref{xxc:eq:theta} is unconditional.
An equivalent $a(n)\sim\tfrac23n^3$ would additionally require
$\Delta(n)=o(n^3)$; the present bounds do not establish that. The computations
are consistent with a much smaller deficit, but this numerical observation
will not be substituted for a nonvanishing theorem.

\subsection{A canonical differential-polynomial representation}

\subsubsection{Three auxiliary variables}
Put
\[
 u=x^x,\qquad v=x^{-1},\qquad z=\log x,\qquad A=z(z+1).
\]
The three elementary derivatives and the logarithmic derivative of $f$ are
\begin{equation}\label{xxc:eq:basicderivatives}
 u'=u(z+1),\qquad v'=-v^2,\qquad z'=v,
 \qquad \frac{f'}{f}=u(A+v).
\end{equation}
There is a unique polynomial $P_n\in\Z[u,v,z]$ such that
\begin{equation}\label{xxc:eq:Pdefinition}
 f^{(n)}(x)=f(x)P_n(x^x,x^{-1},\log x).
\end{equation}
Write
\begin{equation}\label{xxc:eq:coeffdefinition}
 P_n(u,v,z)=\sum_{k,j,\ell\ge0}c_n(k,j,\ell)u^kv^jz^\ell.
\end{equation}
The corresponding expanded term of the original derivative is
\[
 c_n(k,j,\ell)\,x^{x^x+kx-j}(\log x)^\ell.
\]
Therefore
\begin{equation}\label{xxc:eq:count}
 a(n)=\#\{(k,j,\ell):c_n(k,j,\ell)\ne0\}.
\end{equation}
In particular, a negative coefficient contributes one term, not minus one.

\subsubsection{Why the representation is unambiguous}

Suppose a finite linear combination of $u^kv^jz^\ell$, after substitution,
vanishes for all sufficiently large positive $x$. Group it by $k$ and choose
the largest $k$ having a nonzero Laurent-polynomial coefficient in $x$ and
$\log x$. Division by $x^{kx}$ makes every smaller-$k$ group exponentially
small relative to that coefficient. Within the largest group, choose the
smallest $j$, and then the largest $\ell$ with nonzero coefficient. Its
asymptotic magnitude cannot be canceled by the remaining Laurent-logarithmic
terms. This is a contradiction. Hence the displayed monomials are linearly
independent in the required finite sense.

This argument is also a safeguard against a representation-dependent term
count: we are counting a specified canonical basis, not the summands in an
arbitrarily factored expression.

\subsubsection{The recurrence and its proof}

\begin{proposition}[Differential recurrence]\label{xxc:prop:recurrence}
Starting from $P_0=1$, the polynomials satisfy
\begin{equation}\label{xxc:eq:PRec}
 P_{n+1}=u(z^2+z+v)P_n+u(z+1)\partial_uP_n
          -v^2\partial_vP_n+v\partial_zP_n.
\end{equation}
\end{proposition}
\begin{proof}
Differentiate~\eqref{xxc:eq:Pdefinition}, apply the product rule, and substitute
\eqref{xxc:eq:basicderivatives}. The result is exactly~\eqref{xxc:eq:PRec}.
The base case is $f^{(0)}=f$. Induction proves the identity at every order,
and also proves integrality of every coefficient.
\end{proof}

For one coefficient $c=c_n(k,j,\ell)$ the seven contributions are as follows.
All contributions with the same destination must be added before zero
coefficients are removed.
\begin{center}
\begin{tabular}{lll}
\toprule
Destination at order $n+1$ & Contribution & Source in the derivative\\
\midrule
$(k+1,j,\ell+2)$ & $c$ & $f'/f$: $uz^2$\\
$(k+1,j,\ell+1)$ & $c$ & $f'/f$: $uz$\\
$(k+1,j+1,\ell)$ & $c$ & $f'/f$: $uv$\\
$(k,j,\ell+1)$ & $kc$ & differentiation of $u^k$\\
$(k,j,\ell)$ & $kc$ & differentiation of $u^k$\\
$(k,j+1,\ell)$ & $-jc$ & differentiation of $v^j$\\
$(k,j+1,\ell-1)$ & $\ell c$ & differentiation of $z^\ell$\\
\bottomrule
\end{tabular}
\end{center}
Transitions with a zero multiplier are omitted. In particular, no negative
logarithm exponent is ever introduced.

\subsubsection{The first cancellation}

The first two polynomials have the compact forms
\begin{align}
 P_1&=u(A+v),\label{xxc:eq:P1}\\
 P_2&=u^2(A+v)^2+
 u\bigl((z+1)A+(3z+2)v-v^2\bigr).\label{xxc:eq:P2}
\end{align}
Expansion gives $a(1)=3$ and $a(2)=12$. At the next order, the coefficient
polynomial of $uv^2$ is
\[
 -(z+1)+\bigl(\partial_z-1\bigr)(3z+2)=-4z.
\]
Its constant term vanishes. The candidate envelope has $31$ cells at order
$3$, but the derivative has $30$ terms. A Boolean reachability calculation
would fail at this very early example.

\subsection{The support envelope and the cubic upper bound}

\begin{definition}
For $n\ge1$, $1\le k\le n$, and $0\le j\le n$, put
\begin{align}
 m_{n,k,j}&=\max(0,k-j),\label{xxc:eq:lo}\\
 U_{n,k,j}&=
 \begin{cases}
 n+k,&j=0,\\
 \min(n+k-j-1,\,2n-2j),&j\ge1.
 \end{cases}\label{xxc:eq:hi}
\end{align}
Let $\mathcal E_n$ consist of the integer triples with these ranges and
$m_{n,k,j}\le\ell\le U_{n,k,j}$.
\end{definition}

\begin{proposition}[Support containment]\label{xxc:prop:envelope}
For $n\ge1$, $\supp P_n\subseteq\mathcal E_n$.
\end{proposition}
\begin{proof}
Consider a nonzero differentiation path, and let $a,b,c,d,e,r,s$ count its
uses of the seven transitions in the displayed table, in that order. Then
\[
 k=a+b+c,\quad j=c+r+s,\quad
 \ell=2a+b+d-s,\quad
 n=a+b+c+d+e+r+s.
\]
In particular,
\begin{align*}
 n-j&=a+b+d+e,\\
 \ell-(k-j)&=a+d+r\ge0,\\
 2(n-j)-\ell&=b+d+2e+s\ge0,\\
 n+k-j-\ell&=b+c+e+s\ge0.
\end{align*}
If $j>0$, at least one transition of type $c$ or $s$ must occur: the transition
of type $r$ has multiplier $-j$ and cannot create the first positive $j$.
Thus $n+k-j-\ell\ge1$ when $j>0$. The remaining ranges follow immediately
from the transition table and the first differentiation. Every coefficient
is a sum over such paths, so cancellation can remove cells but cannot add a
cell outside these constraints.
\end{proof}

\begin{proposition}[Size of the envelope]\label{xxc:prop:envelopesize}
For $n\ge1$, $\#\mathcal E_n=B(n)$.
\end{proposition}
\begin{proof}
For $j=0$ there are $n+1$ logarithm exponents for each of $n$ values of $k$,
giving $n(n+1)$ cells. For $j\ge1$, set $d=n-j$, so $0\le d\le n-1$.
The number of cells in this $j$-slice is
\begin{align*}
 \sum_{k=1}^n\bigl(\min(d+k,2d+1)-\max(0,k-n+d)\bigr)
 &=n(2d+1)-d(d+1).
\end{align*}
Indeed, replacing each minimum by $2d+1$ overcounts by
$d(d+1)/2$, and the sum of the subtracted maxima is another $d(d+1)/2$.
Consequently
\begin{align*}
 \#\mathcal E_n
 &=n(n+1)+\sum_{d=0}^{n-1}\bigl(n(2d+1)-d(d+1)\bigr)\\
 &=n(n+1)+n^3-\frac{(n-1)n(n+1)}3\\
 &=\frac{2n^3+3n^2+4n}{3}.
\end{align*}
\end{proof}

This proves the upper bound in~\eqref{xxc:eq:mainbounds}. More precisely,
$\Delta(n)$ is exactly the number of zero coefficients \emph{inside}
$\mathcal E_n$. The polynomial $B(n)$ is not being guessed by interpolation:
it is the exact count of a proved finite set of candidate monomials.

\subsection{An independent exact formula for every coefficient}

The recurrence is already an exact algorithm. An independent formula is
valuable for two other reasons: it certifies possible cancellations after a
modular computation, and it proves that every $(k,j)$ row is nonzero.

\subsubsection{Notation for Stirling numbers and auxiliary polynomials}

We use signed Stirling numbers $s(n,k)$ of the first kind and Stirling
numbers $S(n,k)$ of the second kind, with the conventions
\[
 \fall{y}{n}=\prod_{i=0}^{n-1}(y-i)=\sum_k s(n,k)y^k,
 \qquad y^n=\sum_k S(n,k)\fall{y}{k}.
\]
These conventions agree with DLMF \S26.8~\cite{dlmf}. In particular,
$s(n,k)$ has sign $(-1)^{n-k}$ and is nonzero for $1\le k\le n$;
$S(n,k)>0$ on the same range.

For nonnegative $d,p,q$, define the noncentral Stirling numbers
\begin{equation}\label{xxc:eq:Tdef}
 T_d(p;q)=\sum_{r=0}^d\binom dr q^{d-r}S(r,p).
\end{equation}
They can be computed without rational arithmetic:
\begin{equation}\label{xxc:eq:Trec}
 T_0(p;q)=\mathbf1_{p=0},\qquad
 T_{d+1}(p;q)=(p+q)T_d(p;q)+T_d(p-1;q).
\end{equation}
Indices $p<0$ or $p>d$ have value zero. The useful boundary values are
\[
 T_d(0;q)=q^d,\qquad T_d(d;q)=1,\qquad
 T_d(p;0)=S(d,p),\qquad T_d(p;1)=S(d+1,p+1),
\]
where $0^0=1$. Equivalently,
\begin{equation}\label{xxc:eq:Tmoment}
 e^{-w}\sum_{m\ge0}(m+q)^d\frac{w^m}{m!}
 =\sum_{p=0}^dT_d(p;q)w^p.
\end{equation}
To verify this identity, expand $(m+q)^d$, express powers of $m$ in falling
factorials, and sum $\sum_m\fall{m}{p}w^m/m!=w^pe^w$.

Define
\begin{equation}\label{xxc:eq:Qdef}
 Q_{j,r}(\beta)=\fall{\beta}{j}\,\rise{(\beta+1)}{r}
 =\prod_{i=0}^{j-1}(\beta-i)\prod_{i=1}^{r}(\beta+i).
\end{equation}
These integer polynomials are easy to build by multiplying by linear
factors. Notice that $\deg Q_{j,r}=j+r$.

\subsubsection{A parameterized derivative identity}

For parameters $\alpha,\beta$, Taylor expansion with the increment $xt$ gives
\begin{equation}\label{xxc:eq:taylorparam}
 \frac{(x+xt)^{\alpha(x+xt)+\beta}}{x^{\alpha x+\beta}}
 =(1+t)^\beta
 \exp\!\left(\alpha x\,[tz+(1+t)\log(1+t)]\right).
\end{equation}
Comparing coefficients of $t^n$ and then of $(\alpha x)^d$ yields
\begin{equation}\label{xxc:eq:parameteridentity}
 \frac{d^n}{dx^n}x^{\alpha x+\beta}
 =x^{\alpha x+\beta-n}\sum_{d=0}^n(\alpha x)^d
 C_{n,d}(z,\beta),
\end{equation}
where, with $j=n-d$,
\begin{equation}\label{xxc:eq:Cformula}
 C_{n,d}(z,\beta)
 =\sum_{s=0}^d\binom ns\frac{z^s}{(d-s)!}
 \partial_\beta^{\,d-s}Q_{j,d-s}(\beta).
\end{equation}
For completeness, choose $s$ factors $tz$ from the $d$th power of the bracket
in~\eqref{xxc:eq:taylorparam}. The remaining factors contribute
$(1+t)^{d-s}\log^{d-s}(1+t)$. Write the logarithm power as
$\partial_\beta^{d-s}$ acting on $(1+t)^{\beta+d-s}$ and use
\[
 [t^{n-s}](1+t)^{\beta+d-s}
 =\frac{\fall{(\beta+d-s)}{n-s}}{(n-s)!}
 =\frac{Q_{j,d-s}(\beta)}{(n-s)!}.
\]
The factorial prefactors simplify to $\binom ns/(d-s)!$, proving the formula.

\subsubsection{The coefficient oracle}

\begin{theorem}[Exact coefficient formula]\label{xxc:thm:oracle}
Suppose $(k,j,\ell)\in\mathcal E_n$, and put
\[
 d=n-j,\qquad h=d+k-\ell,
\]
\[
 q_- =\max(0,k-\ell,k-d),\qquad q_+=\min(k,h,j).
\]
Then
\begin{equation}\label{xxc:eq:oracle}
\boxed{\displaystyle
 c_n(k,j,\ell)=\sum_{q=q_-}^{q_+}
 \binom n{\ell-k+q}\binom hq\,
 T_d(k-q;q)\,[\beta^h]Q_{j,h-q}(\beta).}
\end{equation}
An empty sum is zero. Outside the envelope, the coefficient is zero.
\end{theorem}
\begin{proof}
Expand the outer exponential:
\[
 f(x)=\sum_{m\ge0}\frac{x^{mx}z^m}{m!},\qquad
 x^{mx}z^m=\left.\partial_\beta^m x^{mx+\beta}\right|_{\beta=0}.
\]
Apply~\eqref{xxc:eq:parameteridentity} with $\alpha=m$. For any polynomial
$C(\beta)$,
\[
 \left.\partial_\beta^m\bigl(e^{\beta z}C(\beta)\bigr)\right|_{\beta=0}
 =\sum_{q=0}^m\frac{m!}{(m-q)!}z^{m-q}[\beta^q]C(\beta).
\]
After division by $f=e^{uz}$, the sum over $m$ occurring for fixed $d,q$ is
\begin{align*}
 e^{-uz}\sum_{m\ge q}\frac{m^d u^m z^{m-q}}{(m-q)!}
 &=u^q e^{-uz}\sum_{r\ge0}(r+q)^d\frac{(uz)^r}{r!}\\
 &=u^q\sum_{p=0}^dT_d(p;q)(uz)^p.
\end{align*}
To obtain $u^kz^\ell$, take $p=k-q$ and take
$s=\ell-k+q$ from~\eqref{xxc:eq:Cformula}. Set $r=d-s=h-q$.
Finally,
\[
 [\beta^q]\frac{\partial_\beta^rQ_{j,r}(\beta)}{r!}
 =\binom{q+r}{q}[\beta^{q+r}]Q_{j,r}(\beta),
\]
and $q+r=h$. This proves~\eqref{xxc:eq:oracle}; the stated bounds on $q$ enforce
$p\in[0,d]$, $s\in[0,d]$, and $q\le j$.

The infinite series can be read as locally normally convergent analytic
series for positive $x$ near any fixed point, or as formal exponential
identities followed by extraction of the finitely many relevant coefficients.
\end{proof}

\subsubsection{Three boundary checks}

Useful independently checkable consequences are
\begin{align}
 [v^0]P_n&=(z+1)^n\sum_{k=1}^nS(n,k)(uz)^k,
                 &&n\ge1,\label{xxc:eq:vzero}\\
 [u^n]P_n&=(z(z+1)+v)^n,\label{xxc:eq:umax}\\
 [v^n]P_n&=\fall{u}{n}.\label{xxc:eq:vmax}
\end{align}
The first follows either from~\eqref{xxc:eq:oracle} or from setting $v=0$ in the
recurrence. The second follows because reaching $u$-degree $n$ requires every
step to differentiate the outer factor $f$. For the third, the coefficient
of the next highest $v$-power is obtained by multiplication by $u-n$,
starting from $1$.

\subsection{A rigorous cubic lower bound}

An upper bound alone does not establish cubic growth. We now show that a
positive cubic number of the candidate cells must survive, without
assuming that cancellations are rare.

\subsubsection{Every row has its advertised highest coefficient}

Let
\[
 R_{n,k,j}(z)=[u^kv^j]P_n(u,v,z).
\]

\begin{lemma}[Nonzero leading coefficient]\label{xxc:lem:leading}
For every $1\le k\le n$ and $0\le j\le n$, the polynomial $R_{n,k,j}$ is
nonzero and has degree exactly $U_{n,k,j}$.
\end{lemma}
\begin{proof}
For $j=0$, equation~\eqref{xxc:eq:vzero} gives
\[
 R_{n,k,0}(z)=S(n,k)z^k(z+1)^n.
\]
Now let $j\ge1$ and put $d=n-j$. If $k\le d$, the proposed highest exponent
is $\ell=d+k-1$, so $h=1$ in~\eqref{xxc:eq:oracle}. Only $q=0,1$ contribute.
The coefficient is
\begin{equation}\label{xxc:eq:leadlowk}
 (-1)^{j-1}(j-1)!\left(
 \binom n{d-1}S(d,k)+\binom nd S(d+1,k)\right),
\end{equation}
which is nonzero. Here we used
$[\beta]Q_{j,0}=[\beta]Q_{j,1}=(-1)^{j-1}(j-1)!$.

If $k\ge d+1$, the highest exponent is $\ell=2d$, and $h=k-d$.
The only possible summand is $q=h$, giving
\begin{equation}\label{xxc:eq:leadhighk}
 [z^{2d}]R_{n,k,j}(z)=\binom nd s(j,k-d).
\end{equation}
It is nonzero because $1\le k-d\le j$. The support envelope already rules
out higher exponents, completing the proof.
\end{proof}

\subsubsection{A high-multiplicity root}

\begin{lemma}[Divisibility]\label{xxc:lem:divisibility}
The row polynomial is divisible by
\begin{equation}\label{xxc:eq:divisibility}
 z^{\max(k-j,0)}(z+1)^{M_{n,k,j}},
 \qquad M_{n,k,j}=\max(n-2j,k-j,0).
\end{equation}
\end{lemma}
\begin{proof}
Extracting $u^kv^j$ from~\eqref{xxc:eq:PRec} gives
\begin{align}\label{xxc:eq:rowrec}
 R_{n+1,k,j}={}&z(z+1)R_{n,k-1,j}+R_{n,k-1,j-1}
              +k(z+1)R_{n,k,j}\notag\\
 &+\bigl(\partial_z-(j-1)\bigr)R_{n,k,j-1}.
\end{align}
Rows with out-of-range indices are zero. Begin with $R_{0,0,0}=1$.
For the factor $z+1$, the first and third terms increase the inherited
multiplicity by one. The second term has inherited multiplicity at least
$\max(n-2j+2,k-j,0)$. The derivative in the last term can reduce its inherited
multiplicity by at most one; before differentiation it is at least
$\max(n-2j+2,k-j+1,0)$. Each is sufficient for the target exponent
$\max(n+1-2j,k-j,0)$. This proves the $(z+1)$ assertion by induction.
The same argument with the factor $z$ proves the exponent $\max(k-j,0)$.
Since $z$ and $z+1$ are relatively prime, both divisibilities hold together.
\end{proof}

\begin{lemma}[Multiplicity forces many terms]\label{xxc:lem:sparse}
A nonzero polynomial with a nonzero root of multiplicity at least $M$ has
at least $M+1$ nonzero coefficients.
\end{lemma}
\begin{proof}
Suppose $p(z)=\sum_{i=1}^t b_i z^{e_i}$, with distinct nonnegative integers
$e_i$ and $b_i\ne0$. If a root $\alpha\ne0$ has multiplicity at least $t$,
then $(z\partial_z)^r p(\alpha)=0$ for $0\le r<t$. Thus
\[
 \sum_{i=1}^t e_i^r b_i\alpha^{e_i}=0,\qquad 0\le r<t.
\]
The matrix $(e_i^r)$ is an invertible Vandermonde matrix, contradicting
$b_i\alpha^{e_i}\ne0$. Therefore the multiplicity is at most $t-1$.
\end{proof}

\subsubsection{Summing the mandatory terms}

By Lemmas~\ref{xxc:lem:leading}--\ref{xxc:lem:sparse}, each $(k,j)$ row contributes at
least $1+M_{n,k,j}$ nonzero monomials. Rows are disjoint in the canonical
basis, so
\begin{equation}\label{xxc:eq:lowersum}
 a(n)\ge\sum_{k=1}^n\sum_{j=0}^n
 \bigl(1+\max(n-2j,k-j,0)\bigr).
\end{equation}
For fixed $j$, the sum of the maxima over $k$ is
\begin{equation}\label{xxc:eq:Mslice}
 \sum_{k=1}^n M_{n,k,j}=
 \begin{cases}
 n(n-2j)+\dfrac{j(j+1)}2,&0\le j\le\lfloor n/2\rfloor,\\[5pt]
 \dfrac{(n-j)(n-j+1)}2,&j>\lfloor n/2\rfloor.
 \end{cases}
\end{equation}
In the first case, start with $n-2j$ in each row and add the excess
$1+\cdots+j$. In the second case, only $k>j$ contributes.
Summing these elementary progressions, and adding the $n(n+1)$ ones
in~\eqref{xxc:eq:lowersum}, gives exactly $L(n)$ in~\eqref{xxc:eq:mainbounds}.
This completes the proof of the growth bounds in Theorem~\ref{xxc:thm:main}.

In particular,
\begin{equation}\label{xxc:eq:limbounds}
 \frac7{24}\le\liminf_{n\to\infty}\frac{a(n)}{n^3}
 \le\limsup_{n\to\infty}\frac{a(n)}{n^3}\le\frac23.
\end{equation}
These constants are proved bounds, not estimates of an experimentally fitted
leading coefficient.

\subsection{Exact cancellation families}

\subsubsection{Logarithm-free coefficients and derivatives of a falling factorial}

Let $F_n(y)=\fall{y}{n}$. In~\eqref{xxc:eq:oracle}, setting $\ell=0$ forces
$q=k$, so for $d=n-j$ and $1\le k\le j$,
\begin{equation}\label{xxc:eq:logfree}
\boxed{\displaystyle
 c_n(k,j,0)=\frac{k^d}{d!\,k!}F_n^{(d+k)}(d).}
\end{equation}
Although written using factorial quotients, this expression is an integer;
the implementation instead uses integer polynomial coefficients and binomial
coefficients. Formula~\eqref{xxc:eq:logfree} reduces every logarithm-free
cancellation to an exact derivative evaluation of a single elementary
polynomial.

\begin{proposition}[Central cancellations at every odd order]\label{xxc:prop:odd}
For odd $n\ge1$,
\begin{equation}\label{xxc:eq:odddeficit}
 \Delta(n)\ge\left\lfloor\frac{n+1}{4}\right\rfloor.
\end{equation}
\end{proposition}
\begin{proof}
Write $n=2d+1$ and set $j=d+1$. Then
\[
 F_n(d+\beta)=\beta\prod_{r=1}^d(\beta^2-r^2)
\]
is odd. For each $1\le k\le d+1$ with $k\equiv d\pmod2$, its coefficient of
$\beta^{d+k}$ vanishes. Formula~\eqref{xxc:eq:logfree} gives a missing monomial
$(k,d+1,0)$. There are $\lfloor(d+1)/2\rfloor=\lfloor(n+1)/4\rfloor$ such
values of $k$.
\end{proof}

Thus $a(n)=B(n)$ is impossible at every odd $n\ge3$. This does not resolve
whether equality occurs infinitely often at even orders. Nor does it
classify all logarithm-free cancellations: additional integer derivative
zeros occur away from the central symmetry.

\subsubsection{An infinite family with positive logarithm exponent}

It is tempting to guess that all cancellations are logarithm-free.
The following family disproves that guess.

Since $f=e^g$ with $g=uz$, the $u^{n-1}$ part of $f^{(n)}/f$ consists of the
partition type having one block of size two and all other blocks of size one.
Equivalently, induction in~\eqref{xxc:eq:PRec} gives
\begin{equation}\label{xxc:eq:almosttop}
 [u^{n-1}]P_n=\binom n2(A+v)^{n-2}
 \bigl((z+1)A+(3z+2)v-v^2\bigr),\qquad n\ge2.
\end{equation}
For $2\le j\le n-1$, put $m=n-j-1$ and
\[
 a=\binom{n-2}{j},\qquad b=\binom{n-2}{j-1},\qquad
 c=\binom{n-2}{j-2}.
\]
The coefficient of $v^j$, after removing the common factor $\binom n2$, is
\begin{equation}\label{xxc:eq:quadraticrow}
 z^m(z+1)^m\left[-cz^2+(a+3b-c)z+(a+2b)\right].
\end{equation}
Its coefficient of $z^{2m+1}$ is
\begin{equation}\label{xxc:eq:nexttotop}
 a+3b-(m+1)c.
\end{equation}
Using
\[
 \frac ac=\frac{m(m+1)}{j(j-1)},\qquad
 \frac bc=\frac{m+1}{j-1},
\]
this vanishes exactly when $m=j^2-4j$ in the stated range.

\begin{proposition}[Non-logarithm-free cancellation family]\label{xxc:prop:family}
For every integer $t\ge2$, let
\begin{equation}\label{xxc:eq:family}
 n=t^2+t-1,\quad k=n-1,\quad j=t+2,\quad \ell=2t^2-7.
\end{equation}
Then $(k,j,\ell)\in\mathcal E_n$, $\ell>0$, and
$c_n(k,j,\ell)=0$.
\end{proposition}
\begin{proof}
Substitute $j=t+2$ into $m=j^2-4j$: it gives $m=t^2-4$ and
$n=m+j+1=t^2+t-1$. Then $\ell=2m+1$ is positive, is inside the envelope,
and its coefficient vanishes by~\eqref{xxc:eq:nexttotop}. The next higher
coefficient is $-\binom n2c\ne0$, so this is a genuine internal missing term.
\end{proof}

The first corresponding orders are $5,11,19,29,41,55,\ldots$, with missing
triples
\[
 (4,4,1),\ (10,5,11),\ (18,6,25),\ (28,7,43),\ldots.
\]
These are additional to the logarithm-free central cancellations. Since all
orders in~\eqref{xxc:eq:family} are odd, this proves
$\Delta(n)\ge\lfloor(n+1)/4\rfloor+1$ along that infinite subsequence.

\subsection{Algorithms and complexity}

\subsubsection{The dependency-free exact recurrence}

The most direct implementation stores a dictionary from triples to integer
coefficients. At each differentiation it applies the seven transitions,
combines equal keys, and removes zeros. A short reference version is given
in Section~\ref{xxc:app:code}. It uses no symbolic-algebra system and no
numerical approximation.

The supplied production version packs a triple into one integer key. For a
requested maximum order $N$, take
\[
 b=N+2,\qquad c=2N+3,\qquad \operatorname{key}(k,j,\ell)=(kb+j)c+\ell.
\]
Incrementing $k,j,\ell$ corresponds to adding $bc,c,1$, respectively. The
chosen bases exceed every possible exponent, so no carries can identify
different monomials. The packed representation reduces tuple allocation
and dictionary overhead; it does not change the mathematics.

\begin{proposition}[Arithmetic complexity]\label{xxc:prop:complexity}
The exact recurrence computes the entire prefix through $N$ in $O(N^4)$
integer arithmetic operations and $O(N^3)$ live integer coefficients.
Each coefficient has $O(N\log N)$ bits.
\end{proposition}
\begin{proof}
At order $n$ there are at most $B(n)=O(n^3)$ terms, and each contributes to
at most seven destinations. Summing through $N$ gives $O(N^4)$ operations.
Only the current and next dictionaries are needed, giving $O(N^3)$ storage.

Let $H_n$ be the sum of the absolute values of the coefficients of $P_n$.
The absolute sum of all multipliers leaving a term is
$3+2k+j+\ell\le3+5n$. Therefore
\[
 H_{n+1}\le(5n+3)H_n,\qquad H_0=1,
 \qquad H_n\le\prod_{r=0}^{n-1}(5r+3)\le5^n n!.
\]
Every individual coefficient satisfies this bound, so its bit length is
$O(n\log n)$.
\end{proof}

This is a unit-cost \emph{integer-arithmetic} bound, not a claim that
arbitrarily large integers are machine words. With elementary small-integer
multiplication, a conservative bit bound is $O(N^5\log^2N)$ time and
$O(N^4\log N)$ coefficient-storage bits. Hashing and object overhead affect
practical memory. We do not claim the algorithm is asymptotically optimal.

\subsubsection{A deterministic modular sieve with exact refinement}

Most coefficients need never be constructed as large integers. Fix any
integer modulus $M\ge2$, and run the same recurrence in $\Z/M\Z$.
If a residue is nonzero, the original integer coefficient is certainly
nonzero. Only zero residues need further attention.

The exact single-term algorithm is:
\begin{enumerate}
\item Compute $P_n\bmod M$ by the seven-transition recurrence.
\item Enumerate the triples in $\mathcal E_n$ whose residues are zero.
\item Evaluate each such coefficient with~\eqref{xxc:eq:oracle}, using integer
      arithmetic, and count only the coefficients that are actually zero.
\item Return $B(n)$ minus that exact zero count.
\end{enumerate}
A zero residue is \emph{never} accepted as proof of a zero coefficient.
Consequently this algorithm is correct for every admissible modulus,
including small or composite moduli. No randomness, primality assumption,
or conjecture about support is required. Tests deliberately use moduli
$2$, $7$, and $12$ to force many false modular zeros.

The supplied NumPy implementation uses signed 64-bit arrays and the modulus
$1\,000\,000\,007$ by default. At order $m$, the array has shape
$(m+1,m+1,2m+1)$. The next array has shape
$(m+2,m+2,2m+3)$, and each transition becomes a shifted slice addition.
All entries are reduced modulo $M$ after each step. The implementation checks
\[
 (10n+10)M<2^{63}
\]
before starting; this conservative condition bounds the intermediate
products and sums before modular reduction. Overflow would invalidate the
certificate, so the check is part of correctness rather than merely an
optimization detail.

The dense modular pass takes $O(n^4)$ fixed-width operations and $O(n^3)$
fixed-width words. Let $Z_M(n)$ denote the number of zero residues inside
the envelope. The oracle tables $Q_{j,r}$ and $T_d(p;q)$, and Pascal's
triangle for binomial coefficients, are built lazily. For a single target
$n$, filling all potentially needed tables takes $O(n^3)$ integer arithmetic
operations and storage; each candidate coefficient then requires at most
$O(n)$ integer products and additions. Thus a conservative oracle bound is
\[
 O(n^3+nZ_M(n))\subseteq O(n^4).
\]
In practice the lazy tables are much smaller when the modular pass leaves
few candidates. Unlike the reference recurrence, the oracle may multiply
two large integers, so its unit-arithmetic bound should again not be read
as a bit-operation bound.

\subsubsection{Which implementation to use}

For a whole prefix, the dependency-free recurrence has the simplest predictable
cost and automatically supplies every intermediate count. For an isolated
larger index, the modular method usually avoids nearly all large-integer
coefficient work. Running a new modular pass separately for every prefix
index would discard this advantage and incur repeated work.

Neither method differentiates a transcendental expression in a computer
algebra system. Both return the exact OEIS term count, including signed
cancellations, rather than an expression-size estimate.

\subsection{Computational validation and results}

\subsubsection{Checks actually performed}

The accompanying files record ten passing test groups. The tests compare the
reference recurrence with all $43$ published A281434 values through $n=42$,
and compare the generated prefix with all $80$ published A352697 deficits
through $n=80$~\cite{oeis281434,oeis352697}. Every candidate coefficient
through $n=12$ is checked against the independent formula~\eqref{xxc:eq:oracle}.
The two methods are also compared through $n=12$ for each of four moduli,
including the deliberately small and composite choices described above.

Additional tests verify the envelope, its exact cardinality, the lower-bound
sum, the $(z+1)$ multiplicities, the two leading-coefficient formulas,
the odd central cancellations, and initial instances of the infinite
non-logarithm-free family. Such finite tests check the implementations and
help detect algebraic mistakes; the universal statements are justified by
the proofs, not by the tests.

The integer recurrence was run through $n=100$. The modular implementation
was separately checked against it at $n=20,40,60,80,100$. Selected larger
computations used the modular pass and exact verification of all candidate
zeros. The archive contains the complete generated prefix, the tests, and
the selected missing triples.

\subsubsection{Selected exact values}

\begin{center}
\begin{tabular}{rrrr}
\toprule
$n$ & $a(n)$ & $B(n)$ & $\Delta(n)$\\
\midrule
0 & 1 & \multicolumn{1}{c}{--} & \multicolumn{1}{c}{--}\\
1 & 3 & 3 & 0\\
2 & 12 & 12 & 0\\
3 & 30 & 31 & 1\\
5 & 113 & 115 & 2\\
11 & 1\,017 & 1\,023 & 6\\
20 & 5\,760 & 5\,760 & 0\\
40 & 44\,320 & 44\,320 & 0\\
60 & 147\,680 & 147\,680 & 0\\
80 & 347\,838 & 347\,840 & 2\\
100 & 676\,800 & 676\,800 & 0\\
101 & 697\,178 & 697\,203 & 25\\
150 & 2\,272\,700 & 2\,272\,700 & 0\\
151 & 2\,318\,265 & 2\,318\,303 & 38\\
200 & 5\,373\,600 & 5\,373\,600 & 0\\
\bottomrule
\end{tabular}
\end{center}
For example, the two missing triples at $n=80$ are
$(33,36,0)$ and $(42,45,0)$. At $n=101$, all $25$ missing terms are supplied
by the central logarithm-free symmetry in Proposition~\ref{xxc:prop:odd}.
These are exact statements about the recorded finite computations.

\subsubsection{Measured performance}

On the execution environment used for this report, the packed integer
recurrence generated the full prefix $0\le n\le100$ in approximately
$26.16$ seconds. Selected single-index modular runs were:
\begin{center}
\begin{tabular}{rrrr}
\toprule
$n$ & Modular zero candidates & Actual zeros & Elapsed seconds\\
\midrule
40 & 0 & 0 & 0.021\\
60 & 0 & 0 & 0.096\\
80 & 2 & 2 & 0.342\\
100 & 0 & 0 & 0.820\\
101 & 25 & 25 & 0.698\\
150 & 0 & 0 & 4.970\\
151 & 38 & 38 & 5.460\\
200 & 0 & 0 & 17.769\\
\bottomrule
\end{tabular}
\end{center}
These are individual measurements, not a controlled benchmark suite or a
promise of timings on other hardware. In particular, computing one target
by the modular method and computing an entire prefix by the integer method
are different workloads. The exact timings, interpreter information, and
modulus are saved in \texttt{data/benchmarks.csv} and
\texttt{data/environment.json}.

The fact that every zero residue in this small selection was a genuine zero
is not assumed by the algorithm. The small-modulus tests exercise the
opposite situation and confirm that false modular zeros are rejected by
the exact oracle.

\subsection{What is and is not established asymptotically}

The unconditional conclusion is cubic-order growth, with the fully explicit
bounds~\eqref{xxc:eq:mainbounds}. Consequently, $a(n)=n^{3+o(1)}$ in the
logarithmic-exponent sense. The lower bound remains valid even if many more
cancellations occur at orders far beyond the computed range.

There are three progressively stronger claims:
\begin{align*}
 \text{Proved here:}\qquad &a(n)=\Theta(n^3).\\
 \text{Not proved here:}\qquad &a(n)\sim\frac23n^3.\\
 \text{Stronger OEIS conjecture:}\qquad
 &a(n)=\frac23n^3+n^2+O(n).
\end{align*}
The last statement is equivalent to $\Delta(n)=O(n)$, given the exact
formula for $B(n)$, and implies the middle statement. Neither follows from
agreement with $B(n)$ at many computed indices.

A proof of the middle statement needs only a subcubic bound on the number
of holes; it need not classify every cancellation. In particular, the
entire plane $\ell=0$ has only $O(n^2)$ candidate cells. Thus even a complete
failure to classify the logarithm-free derivative zeros would not by itself
prevent proving the leading equivalent. The remaining issue is a uniform
bound on cancellations with $\ell>0$ across the three-dimensional envelope.
Formula~\eqref{xxc:eq:oracle} supplies an explicit integer expression for those
coefficients, and the family in Proposition~\ref{xxc:prop:family} shows why a
blanket nonvanishing assertion would be false.

For the stronger linear-deficit conjecture, even the logarithm-free plane
requires more delicate control. Its central symmetry already forces
$\Delta(n)$ to be at least a positive linear function along the odd orders,
so an $o(n)$ bound is impossible. The appropriate target is an upper bound
of matching order, not an assertion that the derivative eventually fills
its whole envelope.

The computation problem is therefore resolved by unconditional exact
algorithms, and the polynomial growth exponent is established. Determining
the exact leading asymptotic constant remains a separate nonvanishing
problem in this analysis.

\subsection{A compact exact implementation}\label{xxc:app:code}

The following version emphasizes transparency. The archive's
\texttt{code/a281434.py} additionally provides packed keys, the coefficient
oracle, modular arrays, input validation, and a command-line interface.

\begin{lstlisting}[language=Python]
from collections import defaultdict

def a281434_prefix(N):
    if not isinstance(N, int) or isinstance(N, bool) or N < 0:
        raise ValueError("N must be a nonnegative integer")
    p = {(0, 0, 0): 1}
    values = [1]
    for _ in range(N):
        q = defaultdict(int)
        for (k, j, ell), c in p.items():
            q[k+1, j, ell+2] += c
            q[k+1, j, ell+1] += c
            q[k+1, j+1, ell] += c
            if k:
                q[k, j, ell+1] += k*c
                q[k, j, ell] += k*c
            if j:
                q[k, j+1, ell] -= j*c
            if ell:
                q[k, j+1, ell-1] += ell*c
        p = {key: c for key, c in q.items() if c}
        values.append(len(p))
    return values
\end{lstlisting}

The only potentially negative transition is the one with multiplier $-j$.
Deleting it, taking absolute values before accumulation, or replacing the
dictionary by a set changes the sequence.

\subsection{Reproducing the supplied computations}

Python 3.9 or later is sufficient for the dependency-free method. The modular
method additionally requires NumPy. From the archive's \texttt{code}
directory, the commands below produce a single term, inspect missing
monomials, create a prefix, and run the tests.
\begin{lstlisting}[language={}]
python a281434.py 100 --method modular
python a281434.py 80 --method modular --holes
python a281434.py 100 --prefix
python -m unittest -v test_a281434.py
python run_experiments.py --reference
python run_experiments.py 20 40 60 80 100 101 150 151 200
\end{lstlisting}
The first command returns \texttt{676800}. The second returns
\texttt{347838} and the two missing triples listed in the article.
The full prefix and related deficits are in \texttt{data/b281434.txt} and
\texttt{data/b352697.txt}. These are generated data files accompanying this
report, not claims that the OEIS has incorporated the additional terms.

The LaTeX source builds with
\begin{lstlisting}[language={}]
pdflatex -interaction=nonstopmode -halt-on-error article.tex
pdflatex -interaction=nonstopmode -halt-on-error article.tex
pdflatex -interaction=nonstopmode -halt-on-error article.tex
\end{lstlisting}
Repeated runs resolve the table of contents and cross-references. The bundled
\texttt{build.sh} and \texttt{build.ps1} perform these commands.

\Needspace{15\baselineskip}

\clearpage
\section{The three answers side by side}\label{syn:sec}

\subsection{What was proved}\label{syn:table}

\begin{center}
\small
\begin{tabular}{@{}p{0.16\textwidth}p{0.25\textwidth}p{0.25\textwidth}p{0.25\textwidth}@{}}
\toprule
 & Part~\ref{xx:part}: $x^{x}$ & Part~\ref{xxb:part}: $x^{x^2}$ &
 Part~\ref{xxc:part}: $x^{x^x}$\\
 & A293239 & A290268 & A281434\\
\midrule
Frame &
 $t=x^{-1}$, $L=\log x$ \newline (2 variables) &
 $z=x^{2}$, $L=\log x$ \newline (2 variables) &
 $u=x^{x}$, $v=x^{-1}$, $z=\log x$ \newline (3 variables)\\[4pt]
Clause of Thm~\ref{fw:thm:canonical} &
 \ref{fw:can:a} &
 \ref{fw:can:b} &
 \ref{fw:can:a}\\[4pt]
Growth order &
 $\Theta(n^{2})$ &
 $\Theta(n^{2})$ &
 $\Theta(n^{3})$\\[4pt]
Upper bound &
 $a(n)\le q(n)$, the proposed value &
 $a(n)\le\U(n)$, the conjectured value &
 $a(n)\le\dfrac{2n^{3}+3n^{2}+4n}{3}$ \newline $(n\ge1$; at $n=0$ it reads
 $0\ge1)$\\[4pt]
Lower bound &
 $n+1+\left\lfloor n^{2}/4\right\rfloor$ &
 $(3n^{2}{+}10n{+}8)/8$ even, \newline $(3n^{2}{+}12n{+}1)/8$ odd &
 $\dfrac{7n^{3}{+}39n^{2}{+}26n{+}3\epsn(n{-}1)}{24}$\\[4pt]
Exact structure &
 $a(n)=1+\binom{n+1}{2}-Z(n)$, \newline $Z$ counting zeros of the
 Lehmer--Comtet triangle A008296 &
 $a(n)=\U(n)-(\text{extra zeros})$, all predicted zeros explained &
 $a(n)=|\EE_n|-\Delta(n)$, $\Delta$ the in-envelope cancellation count\\[4pt]
Settled &
 the reduction; the $\Theta$; a \emph{disproof} of the recurrence printed with
 range $n>6$, at $n=8$ &
 the upper bound; the cancellation mechanism; the $\Theta$; a complete zero
 classification on the first two deficit diagonals &
 the $\Theta$; two exact algorithms; central cancellations at every odd
 order $n\ge3$\\[4pt]
\textbf{Not settled} &
 the principal closed formula: needs Conjecture~\ref{xx:conj:zeros} &
 the conjecture itself, and the constant $\tfrac12$: needs
 Conjecture~\ref{xxb:conj:support} &
 the constant: $a(n)\sim\tfrac23n^{3}$ would need $\Delta(n)=o(n^{3})$\\
\bottomrule
\end{tabular}
\end{center}

\subsection{Why the orders differ}\label{syn:orders}

The growth orders line up with the number of variables in the frame, and
Part~\ref{fw:sec} explains why the upper half of that observation must hold.  In
a frame of $m$ variables, each exponent in $\supp P_n$ is bounded by a constant
multiple of $n$, because each application of the recurrence raises any one
exponent by a bounded amount.  The envelope therefore lies in a box of side
$O(n)$ in $m$ dimensions, so $|\EE_n|=O(n^m)$; with
Corollary~\ref{fw:cor:welldefined} for the two frames it covers, and
Part~\ref{xxc:part}'s own independence argument for the third,
Proposition~\ref{fw:prop:envelope} then gives $a(n)=O(n^m)$.  With $m=2$ for
$x^x$ and $x^{x^2}$ and $m=3$ for $x^{x^x}$, the observed orders $n^2$, $n^2$,
$n^3$ are exactly the envelope dimensions.

Two cautions are worth stating, because the pattern invites more than it
supports.

First, the matching \emph{lower} bounds are not a consequence of the frame at
all.  A frame bounds the support from outside; it says nothing about how much of
the envelope survives cancellation, and in principle almost all of it could
vanish.  Each part proves its own lower bound by exhibiting coefficients that
are nonzero, and those arguments have nothing in common.  Part~\ref{xx:part}
counts nonzero entries along each diagonal of its triangle, bounding the zeros
there by the degree of a polynomial that cannot vanish identically.
Part~\ref{xxb:part} proves strict \emph{positivity} on a region, through a
substitution that turns its coefficients into an integral with a
sign-definite kernel.  Part~\ref{xxc:part} forces a root of high multiplicity
and then uses the fact that a polynomial with a nonzero root of multiplicity
$M$ has at least $M+1$ nonzero coefficients.  Three mechanisms, no common
method; the framework supplies none of them.

Second, the frame is a choice, not an invariant.  Nothing stops one from
embedding any of these problems in a larger frame and obtaining a worse upper
bound; what the three parts exhibit is the smallest frame each author found in
which the derivative closes.  The claim ``the order is the dimension'' is
therefore a statement about these three constructions, not a theorem about power
towers, and it should not be extrapolated to $x^{x^{x^x}}$ without doing the
work.

\subsection{What the shared framework buys, and what it does not}\label{syn:buys}

It buys four things, with two qualifications that should be stated rather than
glossed.  The count is an invariant of the function rather than of a
differentiation algorithm --- by Corollary~\ref{fw:cor:welldefined} for
Parts~\ref{xx:part} and~\ref{xxb:part}, and by Part~\ref{xxc:part}'s own
argument for the third frame, which that corollary does not cover.  The
coefficients are integers, so the computations are exact and a nonzero residue
is a proof of nonvanishing (Proposition~\ref{fw:prop:computation}).  The
\emph{envelope} bound is free once the envelope is described
(Proposition~\ref{fw:prop:envelope}) --- though only Part~\ref{xxc:part}'s
advertised upper bound is the bare envelope count; the bounds listed for the
other two in Section~\ref{syn:table} are the envelope minus proved cancellation
families, and those families are real theorems, not framework.  And the problem
is reduced from ``count the terms'' to ``find the zeros'', which is the form all
three parts actually work in.

It buys nothing at all about the zeros.  That is where each of the three
problems becomes its own problem, and where all three still have something
open.  In
Part~\ref{xx:part} the zeros are the vanishing entries of a classical triangle
and the missing step is their complete classification.  In Part~\ref{xxb:part}
the zeros are forced by a reflection symmetry and the missing step is that there
are no others.  In Part~\ref{xxc:part} the zeros are not even conjecturally
classified, and the open question is the weaker one of their density.  These
three statements are not special cases of one another, and no argument in any
part transfers to any other.

\subsection{What the three do share, beyond the framework}\label{syn:share}

There is more common structure than the framework alone, and it is worth setting
down, because it is the one thing a reader could take from these three parts
that none of them states on its own.

Every cancellation family that any of the three parts \emph{proves} comes from
the same source: the falling factorial $\fall{y}{n}=y(y-1)\cdots(y-n+1)$ is
symmetric about its centre, so that $\fall{(n-1-y)}{n}=(-1)^n\fall{y}{n}$, and a
polynomial that is odd about a point has no even coefficient there.  Each part
reaches that fact in its own coordinates:
\begin{itemize}[leftmargin=1.6em,itemsep=2pt]
\item Part~\ref{xx:part}, Proposition~\ref{xx:prop:symmetry}, centres at $r=2h$
 to get $F_{4h+1}(2h+u)=u\prod_{s=1}^{2h}(u^2-s^2)$, an odd polynomial, killing
 the coefficient of $u^{2h}$;
\item Part~\ref{xxb:part}, Lemma~\ref{xxb:lem:reflection}, uses
 $F_n(n-1-u)=(-1)^nF_n(u)$ directly, transported through a $2$-step difference;
\item Part~\ref{xxc:part}, Proposition~\ref{xxc:prop:odd}, gets
 $F_n(d+\beta)=\beta\prod_{r=1}^d(\beta^2-r^2)$, again odd.
\end{itemize}
So all three counting problems are, at bottom, questions about the integer zeros
of derivatives and finite differences of one classical polynomial, and every
zero any of them can currently explain is a parity zero of that polynomial about
its centre.  What is \emph{not} shared is everything past that point: the
triangles the parity acts on are different, the difference operators are
different, and --- this is the crux --- in no case is the parity family known to
be the whole zero set, which is exactly what each part's open question asks.

Given that, a caution about the two $\Theta(n^2)$ results.  Parts~\ref{xx:part}
and~\ref{xxb:part} both bound a two-dimensional envelope and subtract
cancellations, and both leave the leading constant open, and it would be easy to
read them as one result twice.  They are not: the arrays are different, one is
normalized by $x^{-n}$ and the other is not, and the open statements are
statements about different objects.  What can be said honestly is the
metamathematical thing: no argument given in either part carries over to the
other.  Whether either open conjecture \emph{implies} the other is not something
this article can settle, and it does not claim to.

\subsection{What would settle each one}\label{syn:open}

\begin{itemize}[leftmargin=1.6em,itemsep=3pt]
\item \textbf{$x^x$ (Part~\ref{xx:part}).}  Conjecture~\ref{xx:conj:zeros}: that
 for $1\le r\le m$ the only vanishing first-kind Lehmer--Comtet entries are
 the sporadic $b(8,5)$ and the symmetry family $b(4h+1,2h)$, $h\ge1$.  The
 sporadic entry and the whole family are both \emph{proved} to vanish in
 Part~\ref{xx:part} --- and the sporadic one is genuinely sporadic, not the
 first term of a second family, which is part of what makes the converse hard;
 the open direction is that there are no others.  A proof would give the closed formula for $a(n)$ and,
 with it, the corrected recurrence for the actual derivative count on its true
 range $n\ge14$.
\item \textbf{$x^{x^2}$ (Part~\ref{xxb:part}).}  The nonvanishing statement
 isolated in Conjecture~\ref{xxb:conj:support}: that no coefficient vanishes
 outside the \emph{two} explained families --- the degree-forced top diagonal
 ($j=k$ with $n>2k$) and the reflection family ($k$ even, $n=4k-2j+1$).
 Coefficients do vanish in bulk outside the reflection family alone; it is the
 union that is conjecturally complete.  Theorem~\ref{xxb:thm:equivalence} shows
 the statement is exactly equivalent to the OEIS conjecture, so nothing weaker
 will do.
\item \textbf{$x^{x^x}$ (Part~\ref{xxc:part}).}  A uniform bound
 $\Delta(n)=o(n^3)$ would give the equivalent $a(n)\sim\tfrac23n^3$; the
 stronger $\Delta(n)=O(n)$, conjectured not in A281434 itself but in the
 companion entry A352697 which records the deficits, needs control even on the
 logarithm-free plane, where only $O(n^2)$ cells are candidates.  Whether
 equality $\Delta(n)=0$ recurs infinitely often at even orders is open.
\end{itemize}

The three obstructions do not even have the same shape.  In
Parts~\ref{xx:part} and~\ref{xxb:part} the open statement is a nonvanishing
statement: no coefficient outside the explained families is zero.  In
Part~\ref{xxc:part} it cannot be, because additional zeros provably exist ---
Proposition~\ref{xxc:prop:family} exhibits an infinite family of them with
positive logarithm exponent, and Proposition~\ref{xxc:prop:odd} forces more at
every odd order $n\ge3$.  There the open question is only \emph{how many}: a
density statement, $\Delta(n)=o(n^3)$, not a nonvanishing one.  That difference
is the sharpest thing the comparison has to offer, and it is the reason the
three investigations are printed here as three.

\clearpage
\begin{thebibliography}{99}

\bibitem{oeis293239}
V. Reshetnikov and contributors,
\emph{A293239: Number of terms in the fully expanded $n$th derivative of $x^x$},
The On-Line Encyclopedia of Integer Sequences, 2017.
Conjectured formulas attributed in the entry to C. Barker and V. Kote\v{s}ovec.
Consulted 20 September 2026.
\url{https://oeis.org/A293239}.

\bibitem{oeis290268}
The OEIS Foundation Inc.,
\emph{A290268: Number of terms in the fully expanded $n$-th derivative of
$x^{x^2}$}, contributed by Vladimir Reshetnikov, with formula comments by
Peter Luschny, October 2017. Accessed September 20, 2026.
\url{https://oeis.org/A290268}.

\bibitem{oeis281434}
Vladimir Reshetnikov,
\emph{A281434: Number of terms in the fully expanded $n$-th derivative of
$x^{(x^x)}$}, The On-Line Encyclopedia of Integer Sequences, entry created
October 5, 2017.
\url{https://oeis.org/search?q=id:A281434&fmt=text}.
Accessed September 20, 2026.

\bibitem{oeis8296}
N. J. A. Sloane and contributors,
\emph{A008296: Triangle of Lehmer--Comtet numbers of the first kind},
The On-Line Encyclopedia of Integer Sequences.
Consulted 20 September 2026.
\url{https://oeis.org/A008296}.

\bibitem{oeis352697}
Vladimir Reshetnikov,
\emph{A352697: $a(n)=\mathrm{A037237}(n-1)-\mathrm{A281434}(n)$},
The On-Line Encyclopedia of Integer Sequences, March 29, 2022.
\url{https://oeis.org/A352697}.
Accessed September 20, 2026.

\bibitem{oeis037237}
N. J. A. Sloane and contributors,
\emph{A037237: Expansion of $(3+x^2)/(1-x)^4$},
The On-Line Encyclopedia of Integer Sequences.
\url{https://oeis.org/A037237}.
Accessed September 20, 2026.

\bibitem{comtet}
L. Comtet,
\emph{Advanced Combinatorics: The Art of Finite and Infinite Expansions},
D. Reidel, Dordrecht, 1974, p. 139.

\bibitem{lehmer}
D. H. Lehmer,
\emph{Numbers associated with Stirling numbers and $x^x$},
Rocky Mountain Journal of Mathematics \textbf{15}(2) (1985), 461--479.
\url{https://projecteuclid.org/euclid.rmjm/1250127221}.

\bibitem{gould}
H. W. Gould,
\emph{A set of polynomials associated with the higher derivatives of $y=x^x$},
Rocky Mountain Journal of Mathematics \textbf{26}(2) (1996), 615--625.
\url{https://doi.org/10.1216/rmjm/1181072076}.

\bibitem{jeong}
S. Jeong,
\emph{On a question of Lehmer concerning the Comtet numbers},
arXiv:2607.26613v1, 2026.
In particular, Proposition 3.10(iii) records the symmetry zero family.
The recurrence question answered there concerns the numbers of the
\emph{second} kind, not the complete first-kind zero classification used here.
\url{https://arxiv.org/abs/2607.26613}.

\bibitem{poon}
Steven S. Poon,
\emph{Higher Derivatives of the Falling Factorial and Related
Generalizations of the Stirling and Harmonic Numbers},
arXiv:1401.2737, 2014.
\url{https://arxiv.org/abs/1401.2737}.

\bibitem{dlmf}
NIST Digital Library of Mathematical Functions,
\emph{\S26.8: Set Partitions---Stirling Numbers}.
\url{https://dlmf.nist.gov/26.8}.
Accessed September 20, 2026.

\bibitem{dlmf-reflection}
NIST Digital Library of Mathematical Functions,
\S5.5(ii), especially Eq.~5.5.3, Euler's reflection formula.
Accessed September 20, 2026.
\url{https://dlmf.nist.gov/5.5.E3}.

\bibitem{dlmf-beta}
NIST Digital Library of Mathematical Functions,
\S5.12, especially Eq.~5.12.1, Euler's beta integral.
Accessed September 20, 2026.
\url{https://dlmf.nist.gov/5.12.E1}.

\end{thebibliography}

\end{document}
